\documentclass[11pt]{article}
\usepackage[margin=1.2in]{geometry}
\usepackage{amsmath,amssymb,amsthm}
\usepackage{xr}
\externaldocument{../paper2b_supp/note}
\usepackage{hyperref}
\sloppy

\newtheorem{theorem}{Theorem}
\newtheorem{corollary}[theorem]{Corollary}
\newtheorem{lemma}[theorem]{Lemma}
\newtheorem{conjecture}[theorem]{Conjecture}
\newtheorem{proposition}[theorem]{Proposition}
\theoremstyle{remark}
\newtheorem{remark}[theorem]{Remark}

\newcommand{\E}{\mathbb{E}}
\newcommand{\Q}{\mathbb{Q}}

\title{The biregular case of a spectral-inclusion conjecture:\\
the weighted plane class and its obstructions}
\author{Vico Bonfioli\thanks{Lean~4 sources: \texttt{RamaLean/} in the \texttt{rama-notes} repository. Correspondence: \texttt{vicobonfioli@gmail.com}.}}
\date{August 2026}

\begin{document}
\maketitle

\begin{abstract}
The conjecture that every root of the $d$-matching polynomial $\mu_{d,G}$ lies in
$\operatorname{spec}(T)$, $T$ the universal cover, is false in general \cite{PartI}. Its
biregular case survives, contains Problem~1 of Song, Fan and Miao, and is the subject of this
paper. At $b=2$ the object is the \emph{weighted plane class}: families $\{(c_k,V_k)\}$ of
weighted $2$-planes with $\sum_kc_kP_{V_k}\preceq aI$, whose matching polynomial is the mixed
characteristic polynomial of Marcus, Spielman and Srivastava.

We prove a positive margin for the complete bipartite family, so no $K_{d,q}$ refutes the
repaired statement however wide its gap is made, and we show the margin decays like $n^{-2/3}$,
the soft-edge exponent, so the statement is true but tight and no size-free bound can prove it.
We give an exact vertex recursion for the plane class, decompose its cross-term into pieces whose
nonnegativity would give the band, and prove two unconditional bounds from real-rootedness alone.

The obstructions are the substance. A rank-blind barrier reads the Marchenko--Pastur law and
cannot reach the tree band at either edge; compressions cannot reach the inner end at all; the
coefficient ladder's reach is unbounded in the moment order, but an input at a uniform rate is
the conclusion it would prove, so the dimension restriction is intrinsic; and the tree moment
bound, the natural a priori input, is false off the coordinate case. A ratio route escapes these
and settles the inner edge on a nontrivial class. Everything asserted is machine-checked in
Lean~4/Mathlib unless labelled otherwise.
\end{abstract}

\paragraph{Status conventions.} Claims here are of three kinds, and the text says which each time
rather than leaving it to the reader. A result described as \emph{formalized}, or cited by a Lean
declaration name, is machine-checked in Lean~4/Mathlib: the sources are in \texttt{RamaLean/} of
the accompanying repository, carry no \texttt{sorry}, and rest only on the standard axioms
\texttt{propext}, \texttt{Classical.choice} and \texttt{Quot.sound}. A claim described as
\emph{measured}, \emph{verified to} a stated tolerance, or \emph{checked numerically} is
computational evidence and is not a proof; the script producing it is named, and every named
script is run and its output compared against a recorded baseline
(\texttt{code/cite\_check.py}). A claim called a \emph{conjecture} is unproved, and is called
that even where the evidence is strong.



\section{Preliminaries, from the companion paper}\label{sec:prelim}

This paper is a sequel and does not repeat what it rests on. The conjecture, the closed form at
first Betti number one, the subdivision theorem, the complete bipartite case and the
minimum-degree repair are proved in \cite{PartI}; they are restated here without proof, in the
form used below. Numbering is local to this paper.

\begin{conjecture}\label{conj:gap}
For every finite graph $G$ and every $d\ge1$, all roots of $\mu_{d,G}$ lie in
$\operatorname{spec}(T)$, the spectrum of the universal cover of $G$.
\end{conjecture}

\begin{conjecture}[D3]\label{conj:mindeg3}
Every finite graph of minimum degree at least three satisfies
$\operatorname{Zeros}(\mu_G)\subseteq\operatorname{spec}(T_G)$.
\end{conjecture}

\begin{theorem}\label{thm:main}
For all $n\ge 1$, $r\ge 1$, with $Y=T_n(x/2)$,
\[
\Phi_{n,r}(x)\;=\;U_r(Y)-2U_{r-1}(Y)+U_{r-2}(Y)
\;=\;\underbrace{\bigl(2T_n(x/2)-2\bigr)}_{\chi_{C_n}(x)}\cdot\,
\underbrace{U_{r-1}\bigl(T_n(x/2)\bigr)}_{\Psi_r(x)} .
\]
\end{theorem}

\begin{corollary}\label{cor:gf}
$\displaystyle\sum_{d\ge 0}\mu_{d,C_n}(x)\,z^d=\frac{1}{1-2T_n(x/2)\,z+z^2}$,
\quad and \quad
$\displaystyle\sum_{r\ge 0}\Phi_{n,r}(x)\,z^r=\frac{(1-z)^2}{1-2T_n(x/2)\,z+z^2}$.
\end{corollary}

\begin{corollary}[real-rootedness]\label{cor:roots}
All roots of $\Phi_{n,r}$ lie in $[-2,2]$: with $x=2\cos\theta$,
$\chi_{C_n}$ vanishes at $\theta=2\pi m/n$ and
$\Psi_r=\sin(rn\theta)/\sin(n\theta)$ vanishes at $\theta=m\pi/(rn)$ with $r\nmid m$
(when $r\mid m$ the denominator vanishes too and $\Psi_r=\pm r\neq0$); this leaves
exactly $n(r-1)$ roots, matching $\deg\Psi_r$.
\end{corollary}

\begin{corollary}[Fibonacci--Lucas evaluation]\label{cor:fib}
$\Phi_{n,r}(3)=(L_{2n}-2)\,F_{2nr}/F_{2n}$, since $T_n(3/2)=L_{2n}/2$ and
$U_{r-1}(L_{2n}/2)=F_{2nr}/F_{2n}$.
\end{corollary}

\begin{corollary}[parity]\label{cor:parity}
$\Psi_r(-x)=(-1)^{n(r-1)}\Psi_r(x)$: the quotient is supported on exponents
of a single parity, like a matching polynomial. (For bipartite $C_n$ it is an
even polynomial for every $r$; at $r=2$ it is the classical matching
polynomial $2T_n(x/2)$ of $C_n$, consistent with Godsil--Gutman
\cite{GG}.)
\end{corollary}

\begin{theorem}\label{thm:subdiv}
Let $H$ be $D$-regular, $D\ge2$, and let $S(H)$ be its subdivision, so that the
universal cover of $S(H)$ is the $(D,2)$-biregular tree with spectral gap
$\bigl(0,\sqrt{D-1}-1\bigr)$. Then every nonzero root of $\mu_{S(H)}$ has
absolute value at least $\sqrt{D-1}-1$; that is,
Conjecture~\ref{conj:gap} holds for $S(H)$ at $d=1$. The constant is the spectral
gap of the universal cover; we do not claim it is attained, and by the strictness
of Heilmann--Lieb on finite graphs it is not.
\end{theorem}

\begin{theorem}\label{thm:kpq}
Let $2\le p\le q$. Every root of $\mu_{K_{p,q}}$ lies in
\[
\{0\}\cup\pm\bigl[\sqrt{q-1}-\sqrt{p-1},\ \sqrt{q-1}+\sqrt{p-1}\bigr],
\]
the spectrum of the $(q,p)$-biregular tree. So Conjecture~\ref{conj:gap} holds for
every $K_{p,q}$ at $d=1$.
\end{theorem}

Five displays are used by number below and are reproduced here. The averaging identity
\begin{equation}\label{eq:avg}
\Phi_{n,r}(x)\;=\;\E_{\sigma\in S_r}\;
\prod_{\ell\,\in\,\mathrm{cycles}(\sigma)}\bigl(2T_{n\ell}(x/2)-2\bigr).
\end{equation}
the recurrence carrying the closed form
\begin{equation}\label{eq:V}
V_0=1,\qquad V_1=\mu_G,\qquad
V_d=\mu_G\,V_{d-1}-\mu_{G-V(C)}^{\,2}\,V_{d-2}.
\end{equation}
Newton's identity in the form used for the coefficients
\begin{equation}\label{eq:newton}
r\,\Phi_r=\sum_{k=1}^{r}f(k)\,\Phi_{r-k},\qquad\Phi_0=1 ,
\end{equation}
the arithmetic--geometric mean step
\begin{equation}\label{eq:amgm}
\sqrt{u(u+\alpha)}+1\;\le\;\sqrt{(u+1)(u+1+\alpha)}
\end{equation}
and the two-sided statement
\begin{equation}\label{eq:twosided}
\text{every root of }g\text{ lies in }\bigl[-2\sqrt m,\,2\sqrt m\bigr],
\end{equation}

That is everything borrowed. From here the paper is self-contained, and it proceeds in three
movements: what the biregular case has going for it, what the plane class looks like when the
question is posed operator-theoretically, and what the obstructions are.

\section{Why the biregular case survives, and how far}\label{sec:bireg}

The refutation leaves the biregular case standing, and this section says why and measures how
much room it has: a positive margin for the complete bipartite family, the decay of that margin,
the two obstructions that bar compressions and rank-blind barriers from the inner end, and the
ratio route that escapes them.

\subsection{Why the biregular case survives}\label{sec:biregblock}

For an $(a,b)$-biregular graph the universal cover is the $(a,b)$-biregular tree, whose
spectrum is $\{0\}\cup\{x:\tau\le|x|\le\rho\}$ with
$\tau=|\sqrt{a-1}-\sqrt{b-1}|$ and $\rho=\sqrt{a-1}+\sqrt{b-1}$. Since \cite{HPS} already
places every root in $[-\rho,\rho]$, Conjecture~\ref{conj:gap} for such a graph is
\emph{equivalent} to the assertion that no nonzero root satisfies $|x|<\tau$, which is
Problem~1 of Song, Fan and Miao. That reduction is machine-checked as
\texttt{BiregularBlocking.biregular\_reduction}, with the tree spectrum and the
Hall--Puder--Sawin interval as the two hypotheses; neither is in Mathlib. So the biregular
case is not one fragment of the conjecture among others, it is the whole of what survives.

Both mechanisms of \cite{PartI} are structurally blocked there, and for different
reasons. The subcubic construction hangs branches on a tree skeleton, and a graph in which
every vertex has degree at least two is not a tree, by handshaking against the edge count of
a tree (\texttt{no\_tree\_of\_two\_le\_degree}); a biregular graph with $a,b\ge2$ has minimum
degree at least two and so admits no such skeleton. Hall's construction needs a wide spread of
degrees, and a biregular graph has exactly two. Beyond that, the two constructions place their
roots in \emph{bulk} gaps, at $1.20$, $2.24$ and $2.85$ in the examples above, whereas the
only gap available in the biregular case is the one at the origin: a counterexample there
would need a root just above zero, which is a question about near-deficiency of matchings
rather than about band structure.

There is also a quantitative obstruction, in the direction one would want, though it is not
complete. At a cut vertex with $p$ isomorphic branches the matching polynomial factors through
$F=x\mu_H-p\mu_{H-v}$, and $F(0)=-p\,\mu_{H-v}(0)$, so a Lipschitz bound gives
\[
   |\theta|\;\ge\;\frac{p\,|\mu_{H-v}(0)|}{M},\qquad M=\sup|F'| ,
\]
for every root $\theta$ (\texttt{branch\_root\_sep}). The separation grows linearly in the
number of branches: strengthening the mechanism drives its roots away from zero, which is the
opposite of what a biregular counterexample requires. The bound is conditional on
$\mu_{H-v}(0)\neq0$, that is on $H-v$ having a perfect matching, and Hall's own branch fails
that, since there $H-v$ is a star of odd order. It therefore blocks the subcubic mechanism and
not his, and on the subcubic branch it gives $|\theta|\ge1/6$ against a true smallest root of
$0.662$: valid, and far from sharp.

The vacuity is removable. If $F(0)=0$, divide out the zero: writing $F=x^mG$ with $G(0)\neq0$,
the nonzero roots of $F$ are the roots of $G$, and $G$ is what should be bounded. For Hall's
branch $F=x^5(x^4-16x^2+55)$, so $G(0)=55$ and the obstruction disappears; the bound then
gives $0.991$ against a true root of $\sqrt5$. Carried out on the matching polynomial itself,
where $\mu_G(x)=\pm x^{\,n-2\nu}g(x^2)$ with $\nu$ the matching number and $g(0)=\pm m_\nu$,
the classical Cauchy estimate becomes a statement purely about matching numbers.

\begin{proposition}\label{prop:rootsep}
Let $G$ be a finite graph with matching number $\nu$ and $k$-matching numbers $m_k$. Then
every nonzero root $\theta$ of $\mu_G$ satisfies
\[
   \theta^2\;\ge\;\frac{m_\nu}{m_\nu+\max_{k<\nu}m_k}.
\]
\end{proposition}

The estimate is Cauchy's and is classical; only its reading in terms of matching numbers is
being recorded here. It is machine-checked as \texttt{RootSeparation.cauchy\_lower} and
\texttt{matching\_root\_sep}, with no hypothesis on the graph.

What makes it worth stating is that the inequality points the useful way. By
\texttt{biregular\_reduction}, Conjecture~\ref{conj:gap} for an $(a,b)$-biregular graph is
exactly the absence of a nonzero root below $\tau$, so any graph whose bound reaches $\tau$
satisfies the conjecture, unconditionally and with no appeal to \cite{AFH}; that implication
is \texttt{RootSeparation.sfm\_of\_separation}. Run on random biregular graphs
(\texttt{code/rootsep.py}) it settles three of nine, the smaller ones, the bound falling as
$0.63,0.45,0.28$ against a fixed $\tau=0.318$ for $(3,4)$-biregular as the graph grows. So it
decides small cases outright and is too weak for large ones.

The weakness is not intrinsic: Cauchy's estimate uses only the moduli of the coefficients and
discards the one structural fact available. By Heilmann--Lieb $\mu_G$ is real-rooted, so $g$
has all roots real and \emph{positive}. For a finite family of positive reals the $\ell^m$
norm of the reciprocals decreases to the $\ell^\infty$ norm, so with
$P_m=\sum_i y_i^{-m}$,
\[
   \min_i y_i\;\ge\;P_m^{-1/m},
\]
sharpening as $m$ grows and exact in the limit. The case $m=1$ is the harmonic bound
$m_\nu/m_{\nu-1}$, already better than Cauchy; and the $P_m$ are computed exactly from the
matching numbers by Newton's identities on the reversed polynomial, so every quantity is
rational.

\begin{proposition}\label{prop:powersum}
Let $G$ be biregular with $\tau=|\sqrt{a-1}-\sqrt{b-1}|$, and let $q\ge\tau^2$ be rational.
If $P_m\,q^m\le1$ for some $m\ge1$, then no nonzero root of $\mu_G$ has $|\theta|<\tau$, so
Conjecture~\ref{conj:gap} holds for $G$.
\end{proposition}

\begin{proof}
$y_i^{-m}\le P_m\le q^{-m}$ for each $i$, hence $y_i\ge q\ge\tau^2$.
\end{proof}

Machine-checked as \texttt{PowerSumCertificate.roots\_ge\_of\_powersum}, with division and
$m$-th roots avoided so that the check is exact rational arithmetic. Measured in
\texttt{code/powersum.py}, $m=8$ already agrees with the true smallest root to three or four
decimal places, and the certificate settles all nine biregular graphs tested, including every
size at which the Cauchy and harmonic forms fail.

Run at scale it certifies $417$ of $417$ random biregular graphs across nine degree pairs and
sizes up to $40$ vertices, in exact rational arithmetic throughout
(\texttt{code/certify\_biregular.py}). Each of those is an unconditional proof of
Conjecture~\ref{conj:gap} for that graph.

It is a decision procedure and not a proof of the conjecture, and the distinction is not a
matter of effort. For a finite family of positive reals the $\ell^m$ norm of the reciprocals
\emph{decreases to} the $\ell^\infty$ norm, so
\[
   \sup_m P_m^{-1/m}=y_1
\]
exactly: optimising the certificate over $m$ computes the smallest root rather than bounding
it. Numerically the agreement is already to five decimal places by $m=16$
(\texttt{code/powersum\_limit.py}). Consequently the certificate succeeds for a graph
precisely when $y_1>\tau^2$, which is Song--Fan--Miao for that graph, and a bound on $P_m$
uniform in the size of the graph strong enough to give the conjecture would already be the
conjecture. The difficulty has not moved, and the route is recorded as closed rather than
promising. What the certificate does supply is a complete and exact decision procedure, which
is worth having and is all it is.

Empirically the statement is robust. Over $180$ random $(a,b)$-biregular bipartite graphs on
up to $40$ vertices, for $(a,b)$ ranging over eight pairs, no root falls below $\tau$; the
smallest nonzero root stays a factor $1.25$ to $2.2$ above it
(\texttt{code/biregular.py}). Random sampling is the wrong instrument for a structured
phenomenon, so an adversarial local search was run as well, descending on the smallest nonzero
root by degree-preserving edge swaps, which stay inside the biregular class exactly. It barely
moves the quantity, from $0.6076$ to $0.6005$ over $1200$ swaps in the hardest case, and the
final ratios to $\tau$ run from $1.42$ to $1.97$ (\texttt{code/biregular\_anneal.py}). The
smallest nonzero root appears to be rigid inside the biregular class.

That, together with the first Betti number one case where the conjecture is a theorem, is
where the remaining interest lies.

\subsection{A positive margin for the complete bipartite family}\label{sec:margin}

The biregular case is most exposed at large degree contrast, since the gap $(0,\tau)$ widens
without bound with $q$. For $K_{d,q}$ the question is decidable in closed form. A $k$-matching
picks $k$ vertices on each side and matches them, so
\begin{equation}\label{eq:kdq}
   \mu_{K_{d,q}}=x^{\,q-d}P(x^2),\qquad
   P(y)=\sum_{k\le d}(-1)^k\binom{d}{k}(q)_k\,y^{d-k},
\end{equation}
with $(q)_k$ the falling factorial. Here $P$ is a Laguerre polynomial, and as $q\to\infty$ its
zeros satisfy $(y-q)/\sqrt q\to$ the zeros of the probabilists' Hermite polynomial
$\mathrm{He}_d$. Writing $h_d$ for the largest of those,
$x_{\min}=\sqrt q-\tfrac{h_d}{2}+o(1)$ against $\tau=\sqrt q-\sqrt{d-1}+o(1)$, so the margin
tends to $\sqrt{d-1}-h_d/2$; at $q=10^6$ this agrees to four decimals for every $d$ from $3$ to
$20$ (\texttt{code/d3attack.py}).

\begin{theorem}\label{thm:margin}
For every $d\ge3$ the limiting margin is positive:
$\sqrt{d-1}-h_d/2\ \ge\ \tfrac12\bigl(\sqrt{d-1}-\sqrt{d-2}\bigr)>0$.
\end{theorem}

\begin{proof}
By the recurrence $\mathrm{He}_{n+1}=y\,\mathrm{He}_n-n\,\mathrm{He}_{n-1}$, the zeros of
$\mathrm{He}_d$ are the eigenvalues of the Jacobi matrix with zero diagonal and off-diagonal
entries $\sqrt1,\dots,\sqrt{d-1}$. Its absolute row sums are $\sqrt{i-1}+\sqrt i$, largest in the
last row, so Gershgorin gives $h_d\le\sqrt{d-2}+\sqrt{d-1}$.
\end{proof}

The Gershgorin step is \texttt{CompleteBipartiteMargin.abs\_eigenvalue\_le\_max\_rowsum}, proved
from the eigenvalue equation at an index where the eigenvector is largest; the consequence is
\texttt{margin\_pos\_of\_rowsum}, and \texttt{no\_root\_below\_gap\_edge} closes the chain. No
complete bipartite graph violates Conjecture~\ref{conj:gap}, however wide the gap is made.

\subsection{The margin vanishes, so the statement is tight}\label{sec:softedge}

$K_{d,q}$ is not extremal. Random $(d,q)$-biregular graphs have smaller margins, and the margin
falls as they grow: for $(3,6)$ the minimum over twelve samples falls monotonically across
fourteen sizes, from $0.5301$ at $n=12$ to $0.2009$ at $n=51$, fitting a power law at
$R^2\ge0.999$ in each of twelve families spanning $d=3$ to $6$ (\texttt{code/softedge2.py}).
Matchings are counted exactly, by a bitmask permanent over the smaller side, which is what makes
the range reachable.

The rate is the soft-edge one. The cavity resolvents on the biregular tree satisfy
$X=1/(z-A_1Y)$, $Y=1/(z-B_1X)$ with $A_1=d-1$, $B_1=q-1$, so $P=XY$ solves
$A_1B_1P^2+(A_1+B_1-z^2)P+1=0$; writing $a=\sqrt{A_1}$, $b=\sqrt{B_1}$, its discriminant factors,
\begin{equation}\label{eq:disc}
   (z^2-a^2-b^2)^2-4a^2b^2=\bigl(z^2-(a+b)^2\bigr)\bigl(z^2-(b-a)^2\bigr),
\end{equation}
locating the band edges at $b\pm a$ (\texttt{SoftEdge.discriminant\_factors}). The inner factor
has a simple zero there (\texttt{inner\_edge\_simple}), so the density vanishes like
$\sqrt{x-\tau}$ and the density exponent is $\tfrac12$ for every $(d,q)$. The expected root count
within $\delta$ of the edge is then $CN\delta^{3/2}$, and the extreme root sits where that count
is one, giving $\delta=(CN)^{-2/3}$ (\texttt{edge\_scale}, \texttt{sqrt\_edge\_exponent}).
Writing $\rho(x)=\kappa\sqrt{x-\tau}+O(x-\tau)$, the asymptotic margin is explicit,
\begin{equation}\label{eq:constant}
   \delta=\Bigl(\frac{3}{2\kappa n}\Bigr)^{2/3}
\end{equation}
(\texttt{sqrt\_edge\_quantile}), and agrees with the analytic quantile to four significant
figures at $n=10^6$ in all twelve families. The density itself passes a sharp check:
$n\int_\tau^\rho\rho=r$ to eight decimals, $r$ being the exact number of positive roots.

The consequences are structural, and they disqualify a class of proofs including
Theorem~\ref{thm:margin}. A quantity positive at every size may admit no uniform positive lower
bound, so any argument producing a constant independent of $n$ is insufficient in general
(\texttt{no\_uniform\_lower\_bound}); the Gershgorin constant is such a constant. A lower bound
may not decay more slowly than the true one, so a proof must reach the exponent
(\texttt{exponent\_floor}). And the picture is a Friedman-type one: the edge is the greatest lower
bound of the smallest positive roots and is never attained (\texttt{inf\_not\_attained}), with no
constant above it bounding them below (\texttt{edge\_unimprovable}). There is no $\varepsilon$ to
spare, and the statement cannot be softened.

Two limits on this. The fitted exponents vary with the aspect ratio $q/d$ at the sizes reachable,
from $-0.67$ at $q=2d$ to $-0.61$ at $q=4d$, and the exponent of the analytic quantile also varies,
being $-2/3$ by construction and reaching it only near $n=10^6$. A finite-size effect is the natural
reading, but the check does not establish it: the analytic prediction neither reproduces the measured
dependence nor is flat against it, the exponent gap running to $0.093$, so the script returns
INCONCLUSIVE and the reading stays a reading (\texttt{code/quantile.py}). And the observed smallest roots sit where the tree measure
places about $0.5$ roots, a count that drifts as $n^{+0.10}$; that departure is neither the girth,
which is constant at $4$ throughout, nor the estimator, since the unbiased sample mean drifts
equally, and it is unexplained (\texttt{code/nobs.py}).

\subsection{Compression cannot reach the inner end}\label{sec:barrier}

There is a structural reason the outer end of \eqref{eq:twosided} is a theorem and the inner end
is not. The outer bound is a statement about quadratic forms: every Rayleigh quotient of a
diagonal form lies between its least and greatest entry (\texttt{SoftEdge.rayleigh\_between}), and
being a statement about forms it passes to every subspace. That is exactly the path-tree
argument, the path tree being an induced subtree of $T$ and its adjacency form a compression.

A gap is not a statement about quadratic forms and is not inherited. With entries $-1$ and $1$
the spectrum is $\{-1,1\}$ and $(-1,1)$ is a gap, yet the vector $(1,1)$ has Rayleigh quotient
$0$, strictly inside it (\texttt{rayleigh\_in\_gap}). So a compression can take values the
original excludes, and no argument phrased in Rayleigh quotients or compressions can rule them
out however refined. This is the $P_8$-inside-the-$(3,2)$-biregular-tree phenomenon of
\cite{PartI} at its smallest, and it is why Heilmann--Lieb, Godsil's path tree and the
Marcus--Spielman--Srivastava interlacing machinery all deliver one end of the interval and are
silent about the other.

\subsection{The ratio route}\label{sec:ratio}

What is not barred is the ratio recursion, which bounds ratios of matching polynomials rather
than Rayleigh quotients. By Godsil $\mu_G$ divides $\mu_P$ for the path tree $P$, and on a tree
the deletion recurrence is exact and local, $F_v=\lambda-\sum_u 1/F_u$ over the children $u$ of
$v$ with $F=\lambda$ at leaves, so $\mu_G(\lambda)\ne0$ follows once no $F$ vanishes.

A positive invariant interval cannot do it: invariance would need $a\le\lambda-k/a$, which has a
positive solution only for $\lambda\ge2\sqrt k$, the Heilmann--Lieb bound. The certificate is
instead sign-alternating. On path trees of $(d,q)$-biregular graphs at gap points, over
twenty-two thousand path-tree vertices, the signs separate by vertex type with no exceptions:
every degree-$d$ vertex has $F>0$ and every degree-$q$ vertex has $F<0$ (\texttt{code/pathratio.py}).

\begin{proposition}\label{prop:ratiosteps}
Let $\lambda>0$.
\begin{enumerate}
\item[(i)] If $F_1,\dots,F_k\in[-C,-c]$ with $c>0$ then $\lambda\le\lambda-\sum_j1/F_j\le\lambda+k/c$.
\item[(ii)] If $F_1,\dots,F_k\in[a,B]$ with $0<a\le B$ and $k>\lambda B$, or if
$0<F_j<k/\lambda$ for all $j$, then $\lambda-\sum_j1/F_j<0$.
\end{enumerate}
\end{proposition}

These are \texttt{RatioRoute.left\_step}, \texttt{right\_step} and \texttt{right\_step\_sharp};
(i) holds for every $k$, so a path-tree vertex missing some of its children is still covered, and
the sharp form of (ii) is what the path trees satisfy, with slack $0.42$ to $4.47$ and no
violations over every vertex measured (\texttt{code/childbound.py}).

The combinatorial input is an alternation count. Let $\pi$ end at $w\in R$, put
$k=|N(w)\setminus\pi|$, $p_L=|\pi\cap L|$, $p_R=|\pi\cap R|$, and let $u$ be a child of $w$ with
child count $k_u$.

\begin{lemma}\label{lem:blocking}
$q-r\le k$ and $k_u\le k-(q-r)$.
\end{lemma}

\begin{proof}
The $q-k$ blocked neighbours of $w$ lie in $L$, so $q-k\le p_L$; the path alternates and ends in
$R$, so $p_L\le p_R\le r$, giving $q-r\le k$. For the child, $N(u)\subseteq R$ and an unblocked
neighbour must avoid $\pi\cap R$, so $k_u\le r-p_R\le r-q+k$.
\end{proof}

The three counting steps are formalized in \texttt{PathCount}: \texttt{blocked\_le\_left} and
\texttt{unblocked\_le} are the set-counting steps, and \texttt{alt\_count} is the alternation
count, proved on lists recording which side each path vertex lies on, with
\texttt{alt\_count\_last} transferring it from the first vertex to the last. At the minimum
$k=q-r$ the lemma forces $k_u=0$, so the children are leaves of ratio exactly $\lambda$ and the
parent is $\lambda-k/\lambda<0$ once $k>\lambda^2$.

Assembling both sides gives a system in one constant $m>0$, the lower bound on $|F|$ at
right-type vertices. With $\kappa(t)=\lambda+t/m$ bounding left-type ratios at child count $t$,
and Lemma~\ref{lem:blocking} capping a right-type vertex's children at $k-\Delta$ each where
$\Delta=q-r$, the parent is at most $\lambda-k/\kappa(k-\Delta)$
(\texttt{RatioRoute.right\_step\_bound}); the invariant propagates when that is at most $-m$,
which is
\begin{equation}\label{eq:closure}
   \lambda+m\;\le\;\frac{k}{\kappa(k-\Delta)}\qquad\text{for all }\Delta\le k\le q-1
\end{equation}
(\texttt{closure\_iff}, \texttt{induction\_step}). Equation \eqref{eq:closure} is a fixed-point
condition on $m$ with no closed form; iterating
$m\mapsto\min_k\bigl(k/(\lambda+(k-\Delta)/m)-\lambda\bigr)$ from above gives the largest
admissible $m$.

\begin{theorem}\label{thm:biregclass}
Let $G$ be $(d,q)$-biregular with $r$ vertices of degree $q$, and $\lambda\in(0,\tau)$. If
\eqref{eq:closure} admits some $m>0$, then $\mu_G(\lambda)\neq0$.
\end{theorem}

The recursion consuming the two steps is \texttt{PathTreeInduction.invariant}, by strong induction
on a height function, with \texttt{no\_vanishing} concluding that no ratio vanishes; the path tree
enters only through an abstraction carrying a side map, a children map and a height that strictly
decreases along edges. A single constant is more than the argument needs: the alternation count
applies at every depth, a path of $\ell$ vertices ending in $L$ having $\lfloor\ell/2\rfloor$
vertices in $R$, which caps a left-type vertex's children by
$\min(d-1,\,r-\lfloor\ell/2\rfloor)$ and bounds a right-type vertex's below by
$q-\lfloor\ell/2\rfloor$. Indexing $\kappa$ and $m$ by depth turns \eqref{eq:closure} into a
finite backward recursion, and Theorem~\ref{thm:biregclass} holds verbatim in that form
(\texttt{invariant\_depth}).

Solving numerically, the condition holds throughout the gap for $(d,q,r)$ equal to $(3,6,4)$,
$(3,9,4)$, $(3,12,4)$, $(4,8,4)$, $(4,12,4)$ and $(5,10,4)$, and to $69\%$ of the gap for
$(3,6,5)$ (\texttt{code/depthind.py}); \S\ref{sec:pervertex} removes that last restriction by
bounding each child at its own count rather than at the worst one. This is the first proof of the
inner end of \eqref{eq:twosided} for a nontrivial class.

This has to be reconciled with \S\ref{sec:must}, which rules out cavity recursions. The
obstruction there is that every $b$-regular bipartite graph is an induced subgraph of some
$(a,b)$-biregular graph, and over $b$-regular bipartite graphs the least root tends to zero, so
on the closure of the biregular graphs under vertex deletion the inequality to be proved is
false; an induction carrying the target statement through that closure cannot work.
Theorem~\ref{thm:biregclass} is an induction over vertex-deleted subgraphs, but it does not carry
the target statement through them. It carries the interval bounds, and those are guarded by
$\Delta=q-r\ge1$: when $\Delta\le0$ a right-type path-tree vertex may have no children at all,
the sign alternation fails at once, and \eqref{eq:closure} admits no positive $m$. A $b$-regular
bipartite graph with parts of equal size has $\Delta=b-r<0$, so the family witnessing the
obstruction is exactly what the hypothesis excludes. The condition $\Delta\ge1$ is therefore not
a technical convenience; it is what \S\ref{sec:must} forces, and it is the reason the route
reaches a nontrivial class rather than nothing or everything.

\section{The coefficient ladder}

This also explains the boundary between what is proved and what is not, and the
explanation is exact.

\begin{theorem}\label{thm:subcoef}
Let $G$ be $(a,b)$-biregular bipartite with parts of sizes $p\le q$, degree $a$
on the $p$-side and $b$ on the $q$-side, and let $g(z)=f_G(z+c)$ with
$c=(a-1)+(b-1)$. Then
\[
[z^{p-1}]\,g \;=\; p\,(b-2).
\]
Consequently the roots of $g$ have mean $-(b-2)$, and $g$ is the matching
polynomial of a graph if and only if $b=2$.
\end{theorem}

\begin{proof}
Counting edges by their smaller side gives $pa=qb$, so $p\le q$ forces $a\ge b$
and $b=\min(a,b)$. The polynomial $f_G$ is monic of degree $p$, and its next
coefficient is $-m_1$, where $m_1$ is the number of $1$-matchings, that is
$|E|=pa$. Hence
\[
[z^{p-1}]\,f_G(z+c)=\binom{p}{1}c-pa=p(a+b-2)-pa=p(b-2).
\]
The mean of the roots is $-[z^{p-1}]g/p=-(b-2)$. A matching polynomial of
degree $p$ involves only every other power of $z$, so its coefficient at
$z^{p-1}$ vanishes; thus $b=2$ is necessary. It is sufficient: for $b=2$ the
graph is $S(H)$ for an $a$-regular $H$ and $g=\mu_H$ by \cite{YY}.
\end{proof}

So for $b=2$ the polynomial $g$ \emph{is} a matching polynomial, Heilmann--Lieb
applies verbatim, and one recovers Theorem~\ref{thm:subdiv} and the subdivision
identity as the case $c=(a-1)+1$ of the change of variable. For $b\ge3$ it
provably is not, and Theorem~\ref{thm:subcoef} says more than that it fails: it
locates the failure. The roots of $g$ are not centred at the band centre but sit
a distance $b-2$ below it, towards the lower band edge, which is exactly the
end of \eqref{eq:twosided} that is open. The asymmetry of the problem and the
departure from the matching-polynomial class are the same phenomenon, measured
by the single quantity $b-2$. Examples:
\[
g_{K_{3,4}}=z^3+3z^2-9z-19,\quad
g_{K_{3,5}}=z^3+3z^2-12z-24,\quad
g_{K_{4,5}}=z^4+8z^3-6z^2-128z-139,
\]
with subleading coefficients $3=3(3-2)$, $3=3(3-2)$ and $8=4(4-2)$ as the
theorem requires.

The same method determines more than the one coefficient, and shows that the
graph itself does not enter until surprisingly late.

\begin{theorem}\label{thm:universal}
With notation as in Theorem~\ref{thm:subcoef}, the coefficients of $z^{p}$,
$z^{p-1}$, $z^{p-2}$ and $z^{p-3}$ in $g$ depend only on $(p,a,b)$, not on $G$.
Explicitly the first three are $1$, $p(b-2)$ and
$\tfrac{p}{2}\bigl[(p-1)(b-2)^2-a(b-1)\bigr]$.
\end{theorem}

\begin{proof}
The coefficients of $g$ up to $z^{p-j}$ are determined by $m_1,\dots,m_j$, so it
suffices that $m_1,m_2,m_3$ are functions of $(p,a,b)$ alone. Write $N=|E|=pa=qb$.
Then $m_1=N$. Two distinct edges of a bipartite graph meet in at most one vertex,
so $m_2=\binom N2-p\binom a2-q\binom b2$, which after substituting $q=pa/b$ gives
$m_2=\tfrac{pa}{2}(pa-a-b+1)$; the stated value of $[z^{p-2}]g$ follows from
$\binom p2c^2-(p-1)Nc+m_2$ with $c=a+b-2$.

For $m_3$, let $L$ be the graph on $E$ in which two edges are adjacent when they
meet, so $m_3$ is the number of independent $3$-sets of $L$. Every vertex of $L$
has degree $(a-1)+(b-1)=c$. Three pairwise meeting edges of a bipartite graph
must share a vertex, since otherwise they form a triangle in $G$; hence the
number of triangles of $L$ is $t=p\binom a3+q\binom b3$, and the number of
$2$-edge paths is $N\binom c2-3t$. The standard count
$\binom N3-\tfrac{Nc}{2}(N-2)+\bigl(N\binom c2-3t\bigr)+2t$ therefore depends
only on $(p,a,b)$.
\end{proof}

The graph enters at the very next coefficient, and it does so through a single
invariant, with coefficient one.

\begin{theorem}\label{thm:c4}
Let $C_4(G)$ denote the number of $4$-cycles of $G$. Then
\[
[z^{p-4}]\,g \;=\; U(p,a,b)+C_4(G),
\]
where, writing $\beta=b-2$ and $A=a(b-1)$,
\[
U=\binom p4\beta^4-\tfrac32A\beta^2\binom p3
+\frac{A\,(3A+4\beta^2)}{12}\binom p2-\frac{A\,(9A+2b^2-14b+14)}{24}\,p .
\]
\end{theorem}

\begin{proof}
As in Theorem~\ref{thm:universal}, $m_4$ is the number of independent $4$-sets of
$L$, and inclusion--exclusion expresses it through the \emph{non-induced} counts
in $L$ of the ten graphs on at most four vertices with no isolated vertex. Nine
of the ten are universal, by the clique structure of $L$: each vertex of $L$ is a
cell of the $p\times q$ biadjacency matrix, and two cells are adjacent exactly
when they share a row or share a column, never both, so every edge of $L$ has a
well-defined type.

Only the count of $C_4$ in $L$ is not universal. Reading the cyclic word of types
around a four-cycle of $L$, the words $RRRC$ and $RCCC$ are impossible, since
three consecutive edges of one type force all four cells into a single line and
then the fourth edge has that type too; and $RRCC$ is impossible, since it would
require a cell sharing a column with two distinct cells of one row. This leaves
$RRRR$ and $CCCC$, whose four cells lie in one line and which therefore number
$3\bigl[p\binom a4+q\binom b4\bigr]$, and $RCRC$, whose four cells are exactly
the four entries of a $2\times2$ all-ones submatrix, that is a $4$-cycle of $G$,
each contributing exactly one. Substituting and applying the shift
$[z^{p-4}](z+c)^{p-k}=\binom{p-k}{p-4}c^{4-k}$ gives the stated $U$.
\end{proof}

Together with Theorem~\ref{thm:universal} this determines the top five
coefficients of $g$ completely:
\[
\begin{aligned}
{}[z^{p-1}]g&=p\beta, &
[z^{p-2}]g&=\binom p2\beta^2-\tfrac12Ap,\\
[z^{p-3}]g&=\binom p3\beta^3-A\beta\binom p2+\tfrac16A\beta p, &
[z^{p-4}]g&=U+C_4(G).
\end{aligned}
\]
For the Heawood graph ($p=7$, $a=b=3$, $C_4=0$) the last reads
$35-315+231-77=-126$. So the conjecture is insensitive to the graph until girth
enters, and it enters at the first opportunity: four pairwise disjoint edges are
first constrained by a $4$-cycle.

Theorem~\ref{thm:universal} has a consequence for what a proof of
Conjecture~\ref{conj:gap} can look like. The universal coefficients determine the
first two power sums of the roots of $g$, namely
$\sum\rho_i=-p(b-2)$ and $\sum\rho_i^2=p\bigl[(b-2)^2+a(b-1)\bigr]$, and those
two numbers already imply a bound on the least root.

\begin{proposition}\label{prop:moments}
Every root of $g$ satisfies
$\rho\ge-(b-2)-\sqrt{(p-1)a(b-1)}$, and this is the sharpest bound implied by
the coefficients of $z^p,z^{p-1},z^{p-2}$.
\end{proposition}

\begin{proof}
The roots are real. Fix a root $\rho$ and apply Cauchy--Schwarz to the other
$p-1$: $\bigl(\sum\rho_i-\rho\bigr)^2\le(p-1)\bigl(\sum\rho_i^2-\rho^2\bigr)$.
With $e=\sum\rho_i$ and $S=\sum\rho_i^2$ this is
$p\rho^2-2e\rho+e^2-(p-1)S\le0$, so
$\rho\ge\bigl[e-\sqrt{(p-1)(pS-e^2)}\bigr]/p$. Substituting the two power sums
gives $pS-e^2=p^2a(b-1)$ and the stated bound. Sharpness is the equality case of
Cauchy--Schwarz.
\end{proof}

The bound of Proposition~\ref{prop:moments} grows like $\sqrt p$, whereas the
conjectured bound $-2\sqrt m$ does not depend on $p$ at all. So for all but the
smallest graphs the moment bound is far weaker, and since by
Theorem~\ref{thm:universal} the coefficients it uses do not see $G$, no argument
resting on them can distinguish one $(a,b)$-biregular graph from another. Any
proof of Conjecture~\ref{conj:gap} must therefore use $m_4$ or beyond, that is,
it must be sensitive to the cycle structure of $G$. This is consistent with the
place the conjecture occupies: bounds of the shape $2\sqrt{(a-1)(b-1)}$ are
non-backtracking bounds, and the non-backtracking operator of a bipartite graph
is exactly the object that records cycles. That object can be exhibited, but only
in averaged form, and the reason for the averaging is not incidental.

\section{The multivariate polynomial and the cavity recursion}

The polynomials $g$ admit a clean description. Since a $k$-matching leaves $p-k$
vertices of $A$ unmatched,
\begin{equation}\label{eq:multivar}
f_G(y)=\sum_{k}(-1)^k m_k\,y^{p-k}
=\sum_{M}(-1)^{|M|}\prod_{v\in A\setminus V(M)}y
=\mathcal M\bigl(G;\,x_A=y,\ x_B=1\bigr),
\end{equation}
the multivariate matching polynomial of Heilmann and Lieb under the asymmetric
specialization that puts the variable on one side of the bipartition and $1$ on
the other. So the $g$ are not exotic: they are shifts of matching polynomials in
the multivariate sense. This is a reformulation and not an advance: for
bipartite $G$ every edge of a matching covers one vertex of each side, so
$\mathcal M(G;x_A=y,x_B=z)$ depends only on the product $yz$, and the asymmetric
specialization therefore carries exactly the information of the univariate
matching polynomial. In particular no quantitative Heilmann--Lieb theorem for
non-uniform specializations can help here, and none is available: the
multivariate bounds of \cite{HL} apply a single graph-wide threshold to every
variable.

Deleting a vertex now gives two different recursions, because only $A$-vertices
carry the variable:
\[
v\in A:\quad f_G=y\,f_{G-v}-\!\!\sum_{u\sim v}\!f_{G-v-u},
\qquad
u\in B:\quad f_G=f_{G-u}-\!\!\sum_{v\sim u}\!f_{G-u-v},
\]
whence the cavity recursion
$R_A=y-\sum1/R_B$, $R_B=1-\sum1/R_A$. On the $(a,b)$-biregular tree, with
$P=a-1$, $Q=b-1$, its fixed point satisfies $R^2+(P-Q-y)R+yQ=0$, of
discriminant
\[
D(y)=(P-Q-y)^2-4yQ,
\qquad
D(y)=0\iff y=(\sqrt P\pm\sqrt Q)^2 .
\]
The two roots are exactly the spectrum edges in the variable $y=x^2$, and $D<0$
strictly between them: the fixed point is complex precisely on the spectrum and
real precisely off it, as it must be. Both recursions and the discriminant have
been checked exactly.

Conjecture~\ref{conj:gap} for biregular $G$ is therefore the statement that the
asymmetric recursion, started from the leaves of the path tree, never produces
$R=0$ while $D>0$. The natural invariant region is
\[
R_A\le-\alpha<0,\qquad R_B\in[\,1,\ 1+Q/\alpha\,],
\]
for a parameter $\alpha>0$. A $B$-vertex with children in the region has
$R_B=1+\sum1/|R_A|\in(1,1+Q/\alpha]$, and a $B$-leaf carries $R_B=1$, exactly
the lower boundary. An $A$-vertex with its full complement of $P$ children has
$R_A\le y-P\alpha/(\alpha+Q)$, so the region is closed precisely when
$y\le h(\alpha)$, where $h(\alpha)=\alpha(P-Q-\alpha)/(\alpha+Q)$. The
following is elementary and is machine-checked in Lean~4
(\texttt{region\_closure\_bound}, \texttt{gap\_edge\_identity}):
\begin{equation}\label{eq:square}
(s-t)^2(\alpha+t^2)-\alpha(s^2-t^2-\alpha)=\bigl(\alpha-t(s-t)\bigr)^2 ,
\end{equation}
so $h(\alpha)\le(s-t)^2$ for every $\alpha$, with equality at
$\alpha=t(s-t)=\sqrt{PQ}-Q$. The region therefore closes for every $y$ strictly
below the gap edge and for no $y$ above it: the invariant region is calibrated
exactly to the spectrum, and \eqref{eq:square} is the reason.

\section{The reflection, and why the band edges are never roots}

One structural fact explains why $b=2$ is the case that closes, and it is a
reflection. The subdivision identity
$\mu_{S(H)}(x)=x^{|E|-|V|}\,\mathcal{LM}(H,x^2)$ is due to Yan and Yeh
\cite{YY}, who also treat the regular and semiregular cases; see
\cite{WWM}, Corollary 2.3, for a restatement through the multivariate matching
polynomial. For $H$ $D$-regular it reads $f_G(y)=\mu_H(y-D)$, and
$\mu_H(-z)=(-1)^{|V(H)|}\mu_H(z)$, so
\begin{equation}\label{eq:reflect}
f_G(2c-y)=(-1)^{p}f_G(y),\qquad c=(a-1)+(b-1),
\end{equation}
and $c$ is the midpoint of $[(s-t)^2,(s+t)^2]$ because $(s\pm t)^2=c\pm2\sqrt m$.
The roots of $f_G$ are therefore symmetric about the centre of the band, so for
$b=2$ the lower bound in \eqref{eq:twosided} is \emph{equivalent} to the upper
one, and Theorem~\ref{thm:subdiv} is that equivalence. The reflection
\eqref{eq:reflect} fails as soon as both degrees exceed $2$: it fails for
$K_{3,4}$, $K_{3,5}$ and already for the regular $K_{3,3}$, so regularity does
not rescue it. Since $b=2$ holds exactly when $G$ is the subdivision of an
$a$-regular multigraph, \eqref{eq:reflect} describes the whole of the
phenomenon, and for $b\ge3$ the lower bound is not a repackaging of the upper
bound but carries information the upper bound does not. We note that
\eqref{eq:reflect} has no adjacency analogue: the companion identity
$\varphi(S(H),x)=x^{|E|-|V|}\varphi(H,x^2-D)$ holds as well, but $\varphi(H,\cdot)$
has a parity only for bipartite $H$, whereas $\mu_H$ always does. For the
Petersen graph the matching-side values are symmetric about $D=3$ and the
adjacency-side values are not.

A consequence of the first is that the band edges are never roots, by an
argument of a different kind.

\begin{proposition}\label{prop:galois}
Let $G$ be $(a,b)$-biregular bipartite with $a,b\ge2$, and suppose
$m=(a-1)(b-1)$ is not a perfect square. Then neither $(s-t)^2$ nor $(s+t)^2$ is
a root of $f_G$.
\end{proposition}

\begin{proof}
$f_G$ has integer coefficients and $(s\mp t)^2=(a+b-2)\mp2\sqrt m$, so for $m$
not a square the two band edges are conjugate over $\mathbb Q$; if one is a root
of $f_G$, so is the other. Heilmann and Lieb already give a \emph{strict} bound,
$|a|<2\sqrt{B_1}$ with $B_1=\Delta-1$ for unit edge weights \cite[Thm.~4.3]{HL},
but at the max-degree radius, which for $a\ne b$ exceeds $s+t$; strictness at the
biregular radius therefore has to be argued separately. Every root of $\mu_G$ is
bounded by the spectral
radius of a path tree $T_u(G)$ \cite{Go}, and $T_u(G)$ is finite and embeds
in the universal cover of $G$, the $(a,b)$-biregular tree $T_{a,b}$, with
degrees bounded by those of $T_{a,b}$. Now $\rho(T_{a,b})=s+t$ is the supremum
of $\rho$ over finite subtrees of $T_{a,b}$, and that supremum is not attained:
any finite subtree is a proper subgraph of a larger connected finite subtree,
and $\rho$ is strictly monotone under proper subgraph containment for finite
connected graphs. Hence $\rho(T_u(G))<s+t$, so $(s+t)^2$ is not a root of
$f_G$, and therefore neither is $(s-t)^2$.
\end{proof}

The hypothesis on $m$ is needed for the argument, not merely for convenience:
when $m$ is a square the two band edges are rational and independent, and
conjugacy says nothing. The smallest such case is $(a,b)=(5,2)$, where $m=4$
and the edges are $1$ and $9$. There the conclusion nevertheless holds, but for
a different reason, the Heilmann--Lieb bound is strict for finite graphs, so
Theorem~\ref{thm:subdiv} already excludes the lower edge. Proposition~\ref{prop:galois}
is the statement for the cases those methods do not reach.

Evaluating $f_G$ in $\mathbb Z[\sqrt m]$ confirms this and exhibits the
conjugacy: for $K_{3,4}$, $K_{3,5}$, $K_{3,6}$, $K_{4,5}$ the two edge values
are $53\mp30\sqrt6$, $72\mp40\sqrt8$, $91\mp50\sqrt{10}$, $1877\mp512\sqrt{12}$.
For $b=2$ the irrational parts vanish, as \eqref{eq:reflect} requires.
Proposition~\ref{prop:galois} is worth isolating because it uses no induction on
vertex deletion, and so is not subject to the obstruction of
Section~\ref{sec:must}; it is
the only tool we have found with that property, although it controls only the
edges and not the interior of the gap.

\section{The operator form}

\subsection{The mixed characteristic polynomial}

The multivariate object that does carry the biregular information is the mixed
characteristic polynomial of Marcus, Spielman and Srivastava \cite{MSSII}: with
$P_k=\operatorname{diag}(\mathbf 1_{N(k)})$ of rank $b$ and $\sum_kP_k=aI_p$, one
has $\nu_G=\mu[P_1,\dots,P_q]$, since $\sum_kz_kP_k$ is diagonal and each
$\partial_{z_k}$ must strike a distinct linear factor, so the surviving terms are
the injections, that is the matchings. This is the stability route done properly,
and it locates the shortfall quantitatively: a barrier argument that sees only
the rank, the trace and $\sum_kP_k=aI$ is reading the Marchenko--Pastur law, and
yields the interval $\bigl[(\sqrt a-\sqrt b)^2,(\sqrt a+\sqrt b)^2\bigr]$, which
is strictly wider than the band, the discrepancy being
$2\bigl[\sqrt{ab}-\sqrt{(a-1)(b-1)}-1\bigr]\ge0$ with equality only at $a=b$. So
the standard potentials do not even recover the known upper bound. That is not a
deficiency of the analysis but of the hypotheses, and it can be seen exactly. Take
$A_k=(b/p)I_p$: these are positive semidefinite with
$\operatorname{tr}A_k=b$ and $\sum_kA_k=(qb/p)I=aI$, so they satisfy every
hypothesis the barrier method uses, and $\mu[A_1,\dots,A_q]$ is a Laguerre
polynomial whose roots converge to the Marchenko--Pastur edges. At $(a,b)=(3,2)$ the
upper edge is already exceeded at $p=4$, where the largest root is $6.229$ against
$(s+t)^2=5.828$, and the lower edge at $p=48$, where the smallest is $0.154$
against $(s-t)^2=0.172$. The Marchenko--Pastur interval is therefore attained on
those hypotheses, and \emph{no} barrier argument that does not see
$\operatorname{rank}A_k\le b$ can reach the tree band at either edge. Numerically
the rank condition is what carries the statement: dropping idempotency but keeping
$\operatorname{rank}A_k\le b$ and $\operatorname{tr}A_k=b$ produced no violation in
search, which suggests that class, rather than the projections, as the natural
target. What is
missing is a barrier built from the Kesten measure of the $(a,b)$-biregular tree
rather than from Marchenko--Pastur, which is to say one that uses that the $P_k$
are $0/1$ diagonal, exactly the combinatorial information free probability
discards.

That raises the question of how much of the hypothesis is really needed, and the
answer delimits the problem usefully. Asking for the same conclusion from
\emph{arbitrary} positive semidefinite $A_k$ of rank $b$ with $\sum_kA_k=aI_p$ is
false. Take $p=2$, $q=3$, $(a,b)=(3,2)$, and
$A_1=A_2=\tfrac u2I_2$, $A_3=(3-u)I_2$: each has rank $2=b$ and the sum is $3I=aI$,
while $\mu[A_1,A_2,A_3](y)=y^2-6y+\bigl(6u-\tfrac32u^2\bigr)$. At $u=\tfrac1{10}$
the constant term is $\tfrac{117}{200}$ and the roots are $0.09914$ and $5.90086$,
outside the band $[3-2\sqrt2,3+2\sqrt2]=[0.17157,5.82843]$ at both ends; as
$u\to0$ the roots tend to $0$ and $6$. The counterexample is diagonal, hence
commuting, so noncommutativity is not what fails. What fails is that the traces
are unconstrained: $\operatorname{tr}A_3=5.8\gg b$. Imposing
$\operatorname{tr}A_k=b$ restores the second coefficient, since it then depends on
the family only through $\sum_k\|A_k\|_F^2$, which by Cauchy--Schwarz is at least
$b$ with equality exactly for idempotents. Projections are thus the extreme point
already at the first nontrivial coefficient.

For rank-$b$ orthogonal projections the statement survives, once one adds the
hypothesis $a\ge b$, equivalently $p\le q$, which the graph setting supplies
automatically. Without it the statement is false for a trivial reason: if $p>q$
then $y^{p-q}$ divides $\mu[P_1,\dots,P_q]$, so $0$ is a root while the lower band
edge is positive. Taking $p=6$, $q=4$, $(a,b)=(2,3)$ with $P_1=P_2$ and $P_3=P_4$
the projections onto two complementary $3$-spaces gives a double root at $0$
against a lower edge of $(\sqrt2-1)^2$. A search over more than
$10^5$ feasible families, real and complex projections, and the intermediate
class $\operatorname{rank}\le b$, $\operatorname{tr}=b$, produced no violation,
the smallest observed margin being $3.6\%$ of the band width. That figure needs a
caveat, though, and it cuts against the aggregate impression: every tight case in
those searches has $b=2$, hence lies in the regime settled by
Proposition~\ref{prop:reflect}. In the genuinely open range $b\ge3$ the observed
margins are wide, at $(p,q,a,b)=(6,8,4,3)$ the least root is $0.486$ against a
band edge of $0.101$, a factor of nearly five. The numerical evidence for the open
part of the statement is thinner than the total counts suggest. Two features make
this the right generalization to record. First, commuting rank-$b$ projections with
$\sum_kP_k=aI$ are simultaneously diagonalizable and therefore $0/1$ diagonal, so
they are \emph{exactly} the biregular incidence families: the projection statement
is precisely the noncommutative extension of Conjecture~\ref{conj:gap}. Second, the
universality of Theorem~\ref{thm:universal} persists in that setting, for
projections the coefficients $E_0,\dots,E_3$ depend only on $(p,q,a,b)$, because
idempotency collapses every trace of a word of length at most three with a repeated
letter, while at length four $\operatorname{tr}(A_iA_jA_iA_k)$ does not reduce. The
first free invariant sits at level four on both sides of the analogy, matching the
$4$-cycle count of Theorem~\ref{thm:c4}.

\subsection{A determinantal representation}

The mixed characteristic polynomial admits a probabilistic representation which gives
an unconditional bound with no barrier or stability input. Write $P_k=W_kW_k^{*}$
with $W_k^{*}W_k=I_b$, set $U=[W_1|\cdots|W_q]/\sqrt a$, and let $\Pi=U^{*}U$, a
rank-$p$ orthogonal projection on $\mathbb R^{n}$, $n=pa=qb$, whose $q$ diagonal
$b\times b$ blocks are all $\tfrac1aI_b$. Let $S$ be a sample from the determinantal
process with kernel $\Pi$ and $s_k(S)$ the number of slots of block $k$ it uses.
Then
\begin{equation}\label{eq:dpp}
y^{\,q-p}\,\mu[P_1,\dots,P_q](y)=\mathbb E_S\prod_{k=1}^q\bigl(y-a\,s_k(S)\bigr),
\end{equation}
by expanding $\det(yI+\sum_kz_kP_k)$ over subsets and applying
$\prod_k(1-\partial_{z_k})$, which annihilates every term using a block twice. Here
$|S|=p$ always, so $\sum_ks_k=p$, and each $s_k$ is exactly
$\mathrm{Binomial}(b,1/a)$; for a diagonal family $S$ is simply an independent
uniform choice of one incident edge at each vertex of the $p$-side.

Since $0\le s_k\le b$, every factor in \eqref{eq:dpp} is positive once $y>ab$, so
the expectation is, and there is no root there.

\begin{proposition}\label{prop:abbound}
Every root of $\mu[P_1,\dots,P_q]$ is at most $ab$. Moreover, with
$m=(a-1)(b-1)$,
\[
ab-\bigl(\sqrt{a-1}+\sqrt{b-1}\bigr)^2=\bigl(\sqrt m-1\bigr)^2 .
\]
\end{proposition}

So the bound overshoots the conjectured edge by exactly $(\sqrt m-1)^2$, and is
lossless precisely when $m=1$. At $a=b=2$ the overshoot vanishes
and $ab=4$ is the upper edge; the roots are nonnegative because the matching numbers
are, so \textbf{the band is proved outright at $a=b=2$}.

Where $ab$ stands relative to what is known needs care, and we record it exactly. Xu, Xu
and Zhu \cite[Thm.~1.9]{XXZ} prove, for positive semidefinite $X_1,\dots,X_q$ of rank at
most $k$ with $\sum X_i\preceq I$ and $\varepsilon=\max\operatorname{tr}X_i$, that
$\operatorname{maxroot}\mu\le\bigl(\sqrt{1-\varepsilon/(k-1)}+\sqrt\varepsilon\bigr)^2$
whenever $\varepsilon\le(k-1)^2/k$. Normalizing a rank-$b$ projection family with
$\sum_kP_k=aI$ by $a$ puts it in exactly that form with $k=b$ and $\varepsilon=b/a$, and
since the roots scale with the family this gives
\begin{equation}\label{eq:xxz}
\lambda_{\max}\ \le\ a\Bigl(\sqrt{1-\tfrac{b}{a(b-1)}}+\sqrt{\tfrac ba}\Bigr)^{2},
\qquad a\ \ge\ \frac{b^2}{(b-1)^2},
\end{equation}
which at $b=2$ is $a+2\sqrt{2a-4}$ for $a\ge4$. That is strictly better than $ab$ for
every $a>4$, indeed $a+2\sqrt{2a-4}\le2a$ with equality only at $a=4$, formalized as
\texttt{xxz\_le\_ab} and \texttt{xxz\_eq\_ab\_iff}, and better than
Marchenko--Pastur throughout. So $ab$ is the best available bound only where the
hypothesis of \eqref{eq:xxz} fails, which at $b=2$ means $a<4$: at $a=b=2$, where it is
sharp and settles the band, and at $(a,b)=(3,2)$, where it gives $6$ against a
Marchenko--Pastur edge of $9.899$ and a band edge of $5.828$. Elsewhere \eqref{eq:xxz} is
the state of the art:
\[
\begin{array}{lrrrrr}
\hline
(a,b) & \eqref{eq:xxz} & ab & \text{Marchenko--Pastur} & a+2\sqrt a & \text{band}\\
\hline
(3,2) & \text{n/a} & 6.000 & 9.899 & 6.464 & 5.828\\
(4,2) & 8.000 & 8.000 & 11.657 & 8.000 & 7.464\\
(9,2) & 16.483 & 18.000 & 19.485 & 15.000 & 14.657\\
(100,2) & 128.000 & 200.000 & 130.284 & 120.000 & 119.900\\
(4,3) & 10.977 & 12.000 & 13.928 & - & 9.899\\
(25,4) & 47.126 & 100.000 & 49.000 & - & 43.971\\
\hline
\end{array}
\]
The comparison also says what the target of this section is worth: at $b=2$ the bound
$a+2\sqrt a$ of Proposition~\ref{prop:cross} would improve on \eqref{eq:xxz} for every
$a>4$, since $\sqrt a\le\sqrt{2a-4}$ there (\texttt{target\_le\_xxz}). It is therefore
not a weaker target than the literature already provides, but a stronger one everywhere
the barrier hypothesis holds. Proposition~\ref{prop:abbound} is formalized in
\texttt{RamaLean/ProductBound.lean} as \texttt{no\_root\_above} and
\texttt{edge\_gap}.

The representation also supports an exact recursion on the underlying kernels.
Writing $N_K$ for the transversal expansion attached to a kernel $K$, one has the
border identity
\[
a\,N_{M+ww^{\mathsf T}}=(a-y)\,N_M+N_{\left(\begin{smallmatrix}0&w^{\mathsf T}\\ w&M\end{smallmatrix}\right)},
\]
and Schur complements of admissible kernels remain admissible, with
$\delta_e\le1-c_e$ forced by $K^2\preceq K$. This makes an induction on kernel size
available in principle. It does not close the upper edge: the scalar cavity
invariants that carry the classical argument on graphs are false on the kernel
class, witnessed by explicit rational kernels on which the graph unfolding identity
fails as an inequality by a $y$-independent amount. What survives is a reduction of
the upper edge on the whole class to a single inequality on bordered kernels, proved
for at most two blocks at $b=2$ by monotonicity of transversal determinants under
downdates, which follows from positivity of the adjugate of a positive semidefinite
matrix.

\subsection{The case of rank two}

At $b=2$ the representation \eqref{eq:dpp} collapses further. Only $s_k\in\{0,1,2\}$
occurs, and $\sum_ks_k=p$ forces the counts of $0$s and $1$s to be determined by
$n_2=\#\{k:s_k=2\}$, so with $\Psi(u)=\mathbb E\,u^{n_2}$ one has
\[
\mu(y)=(y-a)^p\,\Psi\bigl(\theta(y)\bigr),\qquad \theta(y)=1-\frac{a^2}{(y-a)^2}.
\]
Real-rootedness of $\mu$ forces $\Psi$ to be real-rooted with nonpositive roots,
hence $n_2$ a sum of independent Bernoulli variables, and the extreme roots of $\mu$
are $a\bigl(1\pm\sqrt{\pi_{\max}}\bigr)$ for $\pi_{\max}$ the largest of those
parameters. Since the band at $b=2$ is $[a-2\sqrt{a-1},\,a+2\sqrt{a-1}]$, both edges
say $a\sqrt{\pi_{\max}}\le2\sqrt{a-1}$:

\begin{proposition}\label{prop:ineq2}
For $b=2$, Conjecture~\ref{conj:gap} for the projection family is equivalent to
\[
\pi_{\max}\ \le\ \frac{4(a-1)}{a^2}\,=\,1-\Bigl(\frac{a-2}{a}\Bigr)^{2}.
\]
\end{proposition}

The equivalence of the two edges with each other and with this inequality is
formalized in \texttt{RamaLean/ProductBound.lean} as \texttt{band\_iff\_pi}; that
they coincide is Proposition~\ref{prop:reflect} seen through $\theta$, which is
two-to-one about $y=a$.

\subsection{The matching form}

Recentring at $a$ turns $\mu$ into a matching polynomial, and the coefficients have
a basis-free form. They also have a Grassmann form: writing
$A_k=b_{k1}b_{k1}^{\mathsf T}+b_{k2}b_{k2}^{\mathsf T}$ and
$\omega_k=b_{k1}\wedge b_{k2}\in\Lambda^2\mathbb R^p$, the Cauchy--Binet formula
gives $e_{2r}\bigl(\sum_{k\in T}A_k\bigr)=\bigl\|\bigwedge_{k\in T}\omega_k\bigr\|^2$,
so
\[
\mu(x+a)=\sum_{T\subseteq[q]}(-1)^{|T|}\Bigl\|\bigwedge_{k\in T}\omega_k\Bigr\|^2x^{\,p-2|T|},
\]
and $M_r$ is a sum of squared Pl\"ucker norms. In the commuting case each norm is
$0$ or $1$ and one recovers the matching numbers.

\begin{theorem}\label{thm:matchingform}
Let $P_1,\dots,P_q$ be rank-two orthogonal projections on $\mathbb R^p$ with
$\sum_kP_k=aI$. Then
\[
\mu(x+a)=\sum_{r\ge0}(-1)^rM_r\,x^{\,p-2r},\qquad
M_r=\sum_{|T|=r}e_{2r}\Bigl(\sum_{k\in T}P_k\Bigr),
\]
where $e_m$ denotes the $m$-th elementary symmetric function of the spectrum. The
$M_r$ are nonnegative, $M_0=1$, $M_1=q$, and $M_r\le\binom qr$.
\end{theorem}

The expression for $M_r$ involves no dilation and no choice of basis for the ranges
of the $P_k$. When the $P_k$ commute, $\sum_{k\in T}P_k$ is a $0/1$ diagonal matrix
of rank $2r$ exactly when $T$ is a matching of the multigraph whose edges are the
blocks, and has rank below $2r$ otherwise; so $e_{2r}$ is the indicator of that
event and $M_r$ is the number of $r$-matchings. Thus at $b=2$ the recentred mixed
characteristic polynomial \emph{is} the matching polynomial in the commuting case,
and $M_r$ is its natural noncommutative extension. Verified exactly on
$S(K_4)$, $S(K_{3,3})$, the cube, and on noncommuting families including the
icosahedral one, where the commutators have norm $0.57$.

Two consequences change how the remaining gap should be described. First, in the
commuting case Proposition~\ref{prop:ineq2} \emph{is} Heilmann--Lieb, so it is a
theorem there, and it is sharp: the extremal families are the large-girth
$a$-regular graphs, for which the ratio of $\pi_{\max}$ to its bound rises
monotonically through $0.866$ at $p=10$ to $0.949$ at $p=38$. No argument with
slack can close it. Second, and more to the point, in that formulation
Proposition~\ref{prop:ineq2} says exactly that the roots of a mixed characteristic
polynomial lie inside the support of the corresponding free convolution, which is
the known open sharpening of the Marcus--Spielman--Srivastava bound. So $b=2$ is
not an easier case awaiting a short argument; it is that open problem in another
notation.

\subsection{The coefficients are mixed discriminants}\label{sec:mixeddisc}

Theorem~\ref{thm:matchingform} writes $M_r$ as a sum of squared Pl\"ucker norms, and
Proposition~\ref{prop:bridge} identifies the coefficients with matching numbers on the commuting
locus. Off that locus they move continuously and remain nonnegative and log-concave. They are not
a new invariant: they are mixed discriminants, the object the mixed characteristic polynomial is
assembled from.

Recall the mixed discriminant of $p$ symmetric $p\times p$ matrices,
\[
D(B_1,\dots,B_p)=\frac1{p!}\sum_{\sigma\in S_p}
\det\bigl[B_{\sigma(1)}\varepsilon_1\ \big|\ \cdots\ \big|\ B_{\sigma(p)}\varepsilon_p\bigr],
\]
introduced by Aleksandrov \cite{Alek} in the theory of mixed volumes. Writing
$\mu(x)=\sum_sc_sx^{\,p-s}$ for a tight rank-$b$ family $A_1,\dots,A_q$ on $\mathbb R^p$,
multilinearity of the determinant gives
\begin{equation}\label{eq:mixeddisc}
c_s=(-1)^s\,s!\,\binom ps\sum_{|S|=s}D\bigl(A_k:k\in S,\ I^{\,p-s}\bigr),
\end{equation}
checked to $1.4\times10^{-14}$ at every $s$ on five families with commutator norms between $0.84$
and $1.00$, so on families far from commuting (\texttt{code/mixeddisc.py}).

Identity~\eqref{eq:mixeddisc} says how to compute the coefficients, not what they are. The following
says what they are, and it needs no new object either.

\begin{proposition}[the coefficients are weighted partial transversals]\label{prop:transversal}
Let $A_1,\dots,A_q$ be rank-$b$ orthogonal projections on $\mathbb R^p$ and let
$u_{k,1},\dots,u_{k,b}$ be an orthonormal basis of $\operatorname{ran}A_k$. Writing
$\mu(y)=\sum_sm_sy^{\,p-s}$,
\begin{equation}\label{eq:transversal}
m_s=(-1)^s\sum_{|S|=s}\ \sum_{r:S\to[b]}\det\operatorname{Gram}\bigl(u_{k,r(k)}:k\in S\bigr).
\end{equation}
\end{proposition}

\begin{proof}
Expanding a determinant perturbed by rank ones,
$\det(yI+\sum_{k,r}w_{k,r}u_{k,r}u_{k,r}^{\mathsf T})
=\sum_Ty^{\,p-|T|}\bigl(\prod_{(k,r)\in T}w_{k,r}\bigr)\det\operatorname{Gram}(T)$,
the sum being over sets $T$ of pairs. Put $w_{k,r}=z_k$, so that the monomial attached to $T$ is
$\prod_kz_k^{t_k}$ with $t_k$ the number of pairs of $T$ in block $k$. Now
$(1-\partial_z)z^t$ evaluated at $z=0$ is $1$ for $t=0$, $-1$ for $t=1$ and $0$ for $t\ge2$, so
applying $\prod_k(1-\partial_{z_k})$ and setting $z=0$ kills every $T$ using two or more vectors from
one block and leaves \eqref{eq:transversal}.
\end{proof}

The inner object is a partial transversal of the block system, one representative chosen from each of
$s$ blocks, weighted by the squared volume they span. On the commuting locus the ranges are
coordinate subspaces, the representatives are standard basis vectors, and the squared volume is $1$
when they are distinct and $0$ when two coincide: the count of partial systems of distinct
representatives, which by Proposition~\ref{prop:bridge} is the matching number of the incidence graph.
Off the locus the weight is the continuous measure of how independent the representatives are. So the
object being weighted off the locus is exactly the object being counted on it. Verified against $\mu$
coefficient by coefficient up to $s=6$ at $C_4$, $K_{3,3}$ and the Fano family, commuting and
deformed (\texttt{code/transversal.py}); no novelty is claimed, the rank-one case being the standard
multi-affine picture of \cite{MSS}. The step the proof turns on, that $(1-\partial_z)z^t$ vanishes at
the origin from $t=2$ on, is machine-checked as
\texttt{Transversal.one\_sub\_deriv\_eval\_zero}, and the two trace identities of
Proposition~\ref{prop:leadrigid} as \texttt{Transversal.sum\_trace\_pairs} and
\texttt{sum\_trace\_triples}, with \texttt{trace\_three\_symm} for the fact that every ordering of a
triple of symmetric blocks carries the same trace. Neither proposition is formalised whole, Mathlib
carrying no mixed characteristic polynomial.

Expanding the Gram determinant makes the rigidity of the leading coefficients a computation rather
than an observation, and in a stronger form than Proposition~\ref{prop:rigid}, which moves along one
rotation.

\begin{proposition}[the four leading coefficients are class-wide constants]\label{prop:leadrigid}
For every tight rank-$b$ projection family at fixed $(n,q,a,b)$,
\[
m_0=1,\quad m_1=-qb,\quad m_2=\binom q2b^2-\tfrac{a^2n-qb}2,\quad
m_3=-\Bigl[\binom q3b^3-\tfrac{b(q-2)(a^2n-qb)}2+\tfrac{a^3n-3a^2n+2qb}3\Bigr].
\]
\end{proposition}

\begin{proof}
The Gram determinant of $s$ unit vectors is $1$ for $s=1$; $1-\langle u,v\rangle^2$ for $s=2$; and
$1-\langle u,v\rangle^2-\langle u,w\rangle^2-\langle v,w\rangle^2
+2\langle u,v\rangle\langle v,w\rangle\langle w,u\rangle$ for $s=3$. Summing over $r$ turns each
closed walk of inner products into a trace, so $m_2$ involves only $\sum_{k<l}\operatorname{tr}A_kA_l$
and $m_3$ only that together with $\sum_{k<l<m}\operatorname{tr}A_kA_lA_m$. Both are fixed by the
hypotheses: from $\sum_kA_k=aI$, $\operatorname{tr}A_k=b$ and $A_k^2=A_k$,
$\operatorname{tr}\bigl((\sum_kA_k)^2\bigr)=a^2n$ gives $\sum_{k<l}\operatorname{tr}A_kA_l
=(a^2n-qb)/2$, and $\operatorname{tr}\bigl((\sum_kA_k)^3\bigr)=a^3n$ gives
$\sum_{k<l<m}\operatorname{tr}A_kA_lA_m=(a^3n-3a^2n+2qb)/6$, every ordered triple of distinct indices
contributing the same value because the blocks are symmetric, so that
$\operatorname{tr}A_kA_mA_l=\operatorname{tr}A_kA_lA_m$.
\end{proof}

At $s=4$ the Gram determinant carries the four-cycle
$\langle u_1,u_2\rangle\langle u_2,u_3\rangle\langle u_3,u_4\rangle\langle u_4,u_1\rangle$, whose sum
over $r$ is $\operatorname{tr}A_kA_lA_mA_r$ with four distinct indices, and the single relation from
$\operatorname{tr}\bigl((\sum_kA_k)^4\bigr)=a^4n$ cannot pin the three independent cyclic classes. So
$m_4$ is the first coefficient carrying a genuine joint invariant, which is why the cut falls at four.
It does move: at the Fano family $m_4$ ranges over an interval of length $6.03$ across random tight
families against the commuting value $1218$, and at $\mathrm{PG}(2,3)$ the four leading coefficients
sit at $1,-52,1170,-14976$ for commuting and deformed families alike while $m_4$, $m_5$ and $m_6$
move from $120627$, $-638820$, $2258100$ (\texttt{code/transversal.py}).

The cycle expansion has a combinatorial reading, and it is one substitution.

\begin{proposition}[the coefficients count systems of distinct representatives]\label{prop:sdr}
For a coordinate family, and for every $S$,
\[
\sum_{\sigma\in\mathrm{Sym}(S)}\mathrm{sgn}(\sigma)\prod_{\text{cycles }c}\Bigl|\bigcap_{k\in c}e_k\Bigr|
=\#\bigl\{f:S\to V \ :\ f(k)\in e_k,\ f\text{ injective}\bigr\},
\]
so $|m_s|$ is the number of partial systems of distinct representatives of size $s$.
\end{proposition}

\begin{proof}
The permutations whose cycle partition is a given $\pi$ contribute
$\prod_{B\in\pi}(-1)^{|B|-1}(|B|-1)!=\mu(\hat0,\pi)$ in total, the M\"obius function of the
partition lattice. The product $\prod_{B\in\pi}|\bigcap_{k\in B}e_k|$ counts the choice functions
constant on every block of $\pi$, so the sum is the M\"obius inversion over the lattice of collisions
and returns the choice functions with no collision.
\end{proof}

This is the classical inclusion-exclusion count of systems of distinct representatives in its
partition-lattice form, and no novelty is claimed for it; it is recorded because of what happens
next. Off the commuting locus the expansion is unchanged except that
\[
\Bigl|\bigcap_{k\in B}e_k\Bigr|\qquad\text{becomes}\qquad \mathrm{tr}\Bigl(\prod_{k\in B}A_k\Bigr).
\]
So $\mu$ is the same signed count for a family of \emph{subspaces}, with the number of points common
to a set of blocks replaced by the cyclic trace of their projections. Three facts recorded separately
above have that single source. The rigidity of $m_0,\dots,m_3$ is that blocks of size at most three
have their cyclic traces fixed by $\sum_kA_k=aI$ and idempotency. The cut at four is that
$\mathrm{tr}(A_iA_jA_kA_l)$ depends on the cyclic order, of which there are three on four blocks,
where an intersection number has no order at all: at a $4$-subset the three orders give one value on
the locus and three off it, spread $0$ against $0.24$ at $K_{3,3}$
(\texttt{code/sdr.py}). And the weights being multi-way rather than pairwise is why no edge-weighted
graph carries them. Verified on $530$ subsets across four families on the locus, and against $\mu$ up
to $s=5$ off it.

The second-order deformation is carried entirely by the four-block data, and that is the structural
statement behind everything above. With $A_k=P_k+\varepsilon D_k+\varepsilon^2X_k$, every cyclic
trace variation is built from the products $D_k(u,v)D_l(u,v)$: two $D$'s give
$\mathrm{tr}(D_xBD_yC)=\sum_{u,v}B(v)C(u)D_x(u,v)D_y(u,v)$ with $B,C$ diagonal, and one $X$ enters
only through its diagonal, which order two forces. So all the variations live in one space, cut by
tightness and the cone condition, and one can ask how much of it each block size spans. The answer is
that nothing beyond four contributes: the rank of the data from blocks of size at most four equals
that from size at most five, six and seven, at $C_6$, $K_{3,3}$, the Fano family and the cube, with
ranks $22$, $82$, $57$ and $189$ (\texttt{code/relations.py}). The relations among the variations are
the cokernel of that map and number $22$, $65$, $28$ and $148$, against the two coming from $P_2$ and
$P_3$ being constant.

That settles the shape of the whole picture. Below four the cyclic traces are fixed by tightness and
idempotency, so $m_0,\dots,m_3$ cannot move. At four the intersection number acquires a cyclic order
and splits into three, which is the only new information the deformation carries. Above four nothing
further enters, so $\delta m_5$ and $\delta m_4$ are two functionals on the same space and can be
proportional, while $\delta m_6$ is a third functional on it and need not be.

Half of the ceiling is a counting argument rather than a computation.

\begin{proposition}[long cycles do not contribute]\label{prop:arcbound}
The second-order row of a cyclic word of length $m$ vanishes identically whenever $m>2a$.
\end{proposition}

\begin{proof}
A nonzero coefficient at $z^{uv}_{xy}$ requires $u\in I_1$ and $v\in I_2$, or the same with $u,v$
exchanged. Every block on the first arc then contains $u$, and $x$ and $y$ both split the pair
$\{u,v\}$, so whichever of them contains $u$ is a further block containing $u$ and not on that arc.
A point of a tight family lies in exactly $a$ blocks, so the first arc has at most $a-1$ of them, and
the second at most $a-1$ likewise; if instead $x$ and $y$ both contain $u$ the first arc has at most
$a-2$ and the second at most $a$. Either way $m=2+|\text{arc}_1|+|\text{arc}_2|\le2a$. The diagonal
terms need every other block to contain the same point and so require $m\le a+1$.
\end{proof}

The counting step is machine-checked as \texttt{ArcBound.card\_le\_of\_mem\_all} and
\texttt{arc\_bound}. The bound is tight: at $K_{3,3}$ and the cube, both with $a=3$, rows of length
exactly $6=2a$ are nonzero (\texttt{code/ceiling.py}). At $a=2$ it reads $m\le4$ and closes the
four-block ceiling outright, with no linear algebra. At $a\ge3$ it leaves rows of length five and six
nonzero, and their dependence on the shorter ones, which the ranks above exhibit, is not explained
here. That dependence is not local either: a length-five row lies in the span of the rows on its own
five blocks at $K_{3,3}$ and at the cube, and does not at the Fano family, so no single local identity
accounts for it and a proof for $a\ge3$ has to be global.

At the level of these rows the second half of Conjecture~\ref{conj:lockstep} is not conjectural. It
follows from the shape of a single difference, and the mechanism is that the difference lies in the
ideal of functionals that tightness makes identically zero.

\begin{proposition}[the five-block row is determined by the four-block row]\label{prop:I2}
Write $W_4^{\mathrm{row}}$ and $W_5^{\mathrm{row}}$ for the second-order rows of $W_4$ and $W_5$ in
the coordinates $z^{uv}_{kl}=D_k(u,v)D_l(u,v)$, and suppose
\[
W_5^{\mathrm{row}}-4(a-2)\,W_4^{\mathrm{row}}=c\,T,\qquad
T(z)=\sum_{u<v}\ \sum_{k,l\in K(u,v)}z^{uv}_{kl},
\]
for some scalar $c$. Then $\delta W_5=4(a-2)\,\delta W_4$ for every direction and every curve.
\end{proposition}

\begin{proof}
$T$ is the vanishing functional of tightness. Indeed
\[
T(z)=\sum_{u<v}\ \sum_{k,l}D_k(u,v)D_l(u,v)=\sum_{u<v}\Bigl(\sum_kD_k(u,v)\Bigr)^2,
\]
and $\sum_kD_k=0$ is the linearisation of $\sum_kA_k=aI$, so $T$ vanishes at every admissible $z$.
Hence $W_5^{\mathrm{row}}$ and $4(a-2)W_4^{\mathrm{row}}$ agree as functionals on the accessible data,
which is the only place either is evaluated.
\end{proof}

The hypothesis is a statement about coefficient patterns and is checked directly: the difference takes
exactly one value on the diagonal coordinates $z^{uv}_{kk}$ and exactly twice that on the off-diagonal
$z^{uv}_{kl}$, the doubling being the $k<l$ storage convention, at $C_6$, $C_8$, $K_{3,3}$, the Fano
family and the cube (\texttt{code/ceiling.py}). The scalar is not derived. It is $0$, $-12$ and $-60$
at $a=2,3,4$, the same at $q=7$, $9$ and $12$ and so independent of $q$; a prediction of $-12(a-2)$
frozen before the $a=4$ computation missed, giving $-24$ against $-60$. None of that matters for the
proposition, whose conclusion holds whatever $c$ is.

The first half of the conjecture, $\delta R=2(a-2)\delta Q_2$, has no such argument. It holds exactly,
to between $10^{-15}$ and $4\times10^{-14}$ at all four families, as a certificate: the functional
$\phi=(d_P-2(a-2)e_P;\,c_T)$ annihilates the image of the $z$-map. But its difference is not a
multiple of $T$. That shape holds at the Fano family and $K_{3,3}$, with $c=-6$ and $-3$ respectively,
and fails at $C_6$ and the cube, where the difference takes several values on each kind of coordinate;
the weaker form $\sum_{u<v}\alpha_{uv}T^{uv}$, one coefficient per vertex pair, fails at the same two.
So the first half must also engage the quadrics $Q_j$, and it rests on a per-family certificate.
Conjecture~\ref{conj:lockstep} is therefore stated as a conjecture, with its second half proved.

Among the coefficients that do move, the first two are not independent, and this is the sharpest
statement the class supports.

\begin{proposition}[the second-order variation of $m_5$]\label{prop:lockstep}
Let $A_0$ be a commuting tight rank-$b$ projection family and $A(\varepsilon)$ a curve in the variety
with $A(0)=A_0$. Every coefficient of $\mu$ varies to order $\varepsilon^2$, the first order
vanishing because the blocks are diagonal and the tangent directions off-diagonal, and
\[
\delta m_5=-b(q-4)\,\delta m_4+2\bigl(\delta R-\delta W_5\bigr),
\]
with $R$ the sum over disjoint triple-and-pair of $\tau\,t$ and $W_5$ the five-cycle trace sum.
\end{proposition}

\begin{conjecture}[the first two moving coefficients are locked]\label{conj:lockstep}
$\delta R-\delta W_5=2(a-2)\,\delta m_4$, so that $\delta m_5=c\,\delta m_4$ with
$c=4(a-2)-b(q-4)$ an integer independent of the direction and of the curve.
\end{conjecture}

Proposition~\ref{prop:lockstep} is the derived part. Expanding the cycle formula at $s=5$ and
summing over $5$-subsets,
\[
m_5=-\Bigl[\tbinom q5b^5-b^3\tbinom{q-2}3P_2+b(q-4)Q_2+2b^2\tbinom{q-3}2P_3-2R-2b(q-4)W_4+2W_5\Bigr],
\]
with $R$ the sum over disjoint triple-and-pair of $\tau\,t$ and $W_5$ the five-cycle trace sum. Since
$P_2$ and $P_3$ are class-wide constants and $\delta m_4=\delta Q_2-2\delta W_4$,
\[
\delta m_5=-b(q-4)\,\delta m_4+2\bigl(\delta R-\delta W_5\bigr).
\]
which is the proposition. A universal $c$ therefore requires $\delta R-\delta W_5$ to be a universal
multiple of $\delta m_4$, which is Conjecture~\ref{conj:lockstep} and is measured rather than proved.
The evidence: the ratio is direction-independent to between $10^{-9}$ and $10^{-7}$ at six
families, and unchanged when the second-order term of the curve is perturbed at scales $1$ and $3$,
so it is a property of the variety and not of the retraction. The formula was fitted on $C_6$, $C_8$,
$K_{3,3}$ and the Fano family, giving $-4$, $-8$, $-6$, $-5$, and then predicted $0$ for the
$2$-regular $3$-uniform family on six vertices and $-16$ for $\mathrm{AG}(2,3)$ before those were
computed; both were hit, at $0.00000000$ and $-15.99999999$ (\texttt{code/lockstep.py}).

The lock does not extend. The ratio $\delta m_6/\delta m_4$ varies over directions by $1.6$ at $C_6$
and $3.6$ at $K_{3,3}$, so the second-order deformation is constrained but not one-dimensional. Taken
with Proposition~\ref{prop:leadrigid}, the second-order variation of $\mu$ lies in a subspace of
codimension at least five in the $(n+1)$-dimensional coefficient space, for every family and every
direction.

One reading is excluded outright. A weighted matching polynomial would have the transversal weight
factor through pairwise data, and it does not, for a reason visible on the commuting locus with no
deformation. In the Fano plane any two lines meet in one point, so every triple of blocks carries the
same pairwise traces; but three lines are either concurrent or form a triangle, and
$\operatorname{tr}P_kP_lP_m$ is $1$ or $0$ accordingly. Two triples with identical pairwise weights
and different transversal weights is exactly the failure of pairwise factorisation, and the same
happens at $\mathrm{PG}(2,3)$. The weight is irreducibly higher than pairwise, and no edge-weighted
graph carries it.

Two properties of $D$ have to hold before \eqref{eq:mixeddisc} even parses, and both are now
machine-checked rather than carried as prose. The sum on the right is over \emph{sets} $S$, which is
legitimate only because $D$ is symmetric in its arguments
(\texttt{MixedDiscriminant.mixedDisc\_comp\_perm}); and calling $D$ the multilinear extension of the
determinant is a claim whose two halves are $D(B,\dots,B)=\det B$
(\texttt{mixedDisc\_const}) and additivity in each argument separately
(\texttt{mixedDisc\_update\_add}), the latter holding because in the term indexed by $\sigma$ the
matrix $B_k$ occupies exactly one column, the one numbered $\sigma^{-1}(k)$. Aleksandrov's
inequality, Bapat's theorem and Gurvits' theorem are cited and not formalised, and neither is
\eqref{eq:mixeddisc} itself.

Two properties recorded earlier as observations are then consequences of classical results rather
than facts about this construction. The $c_s$ alternate in sign because a mixed discriminant of
positive semidefinite arguments is nonnegative, which is Aleksandrov's inequality; the
quantitative theory of that quantity is Bapat's \cite{Bapat}, whose conjectured sharp lower bound
for doubly stochastic tuples was proved by Gurvits \cite{Gurvits}. And the coefficients are
log-concave because $\mu$ is real-rooted, by \cite{MSS}, so Newton's inequalities apply; the
measured ratios run from $1.07$ to $1.20$.

The same identification places the note's own obstruction in classical territory. A tight family
with $q=p$ and $a=b$, scaled by $1/a$, has $\sum_kA_k=I$ and $\operatorname{tr}A_k=1$, so it is a
doubly stochastic tuple in the sense of \cite{Gurvits}, and its mixed discriminant is at least
$p!/p^p$ with equality only at $A_k=I/p$. That tuple is exactly the one recorded in
\S\ref{sec:leaves} as satisfying rank, trace and $\sum_kA_k=aI$ while exceeding the band: its
greatest root is $4.698$ at $p=4$ against a band edge of $4$, and it grows with $p$. The example
that blocks every argument reading only those three data is therefore the van der Waerden extremal,
not an ad hoc construction. Both halves are checked in \texttt{code/mixeddisc.py}, in the
normalisation with $D(I,\dots,I)=p!$.

What \eqref{eq:mixeddisc} settles is the question of what the coefficients count away from the
commuting locus. On the locus, the mixed discriminant of coordinate projections is a count, of
systems of distinct representatives, which is Proposition~\ref{prop:bridge}. Off it, it is the
multilinear extension of that count and is not itself a count. No novelty is claimed for
\eqref{eq:mixeddisc}; it is recorded because the alternative was to keep describing a classical
object in ad hoc terms.

\subsection{The plane class, and two unconditional bounds}

This subsection is the longest in the note, so its outcome is stated first. It delivers an exact
vertex recursion for the plane class (Theorem~\ref{thm:vertexrec}), checked numerically to
$10^{-13}$; a decomposition of its cross-term into pieces $C_r$ (Proposition~\ref{prop:cross},
Theorems~\ref{thm:C2} and~\ref{thm:Cr}) whose nonnegativity would give the band
$[a-2\sqrt a,\,a+2\sqrt a]$ for all weighted plane families, supported numerically but
conjectural; and two bounds following unconditionally from real-rootedness alone, giving the band
whenever $m\le8$ and whenever $m\le16+\bigl(\sum_kc_k^2\bigr)/a^2$, the latter reproducing the
known two-moment bound for projections and extending it to all weighted plane families. The
general case is not closed. A reader who wants only that summary may skip to \S\ref{sec:rank}.

What the problem is, at $b=2$, can be said exactly. Both $\mu(x+a)$ and the operator
$\operatorname{Adj}(A):=\sum_k e_2(A_k)P_{\operatorname{range}A_k}$ depend on the
family only through the pairs $\bigl(e_2(A_k),\operatorname{range}A_k\bigr)$. So the
class is the class of weighted $2$-plane families
$\{(c_k,V_k):c_k\ge0,\ \sum_kc_kP_{V_k}\preceq aI\}$, and $\mu(x+a)$ is the matching
polynomial of such a family; for coordinate planes it is the ordinary weighted
matching polynomial. Two points about that identification need care, and both are invisible in
the projection case.

First, the invariance holds within the tight class and not outside it: with $\sum_kA_k$ a
multiple of the identity, families sharing the pairs
$\bigl(e_2(A_k),\operatorname{range}A_k\bigr)$ have the same recentred $\mu$ even when the $A_k$
are strongly anisotropic, verified to $10^{-13}$ at eigenvalue ratios up to $129$
(\texttt{code/rl\_conj.py}); without tightness the recentring has no meaning.

Second, the recentring constant is that tight constant, and it is \emph{not} the constant in
$\operatorname{Adj}(A)\preceq aI$. Plane data with $\operatorname{Adj}(A)=aI$ admits a
realization with $\sum_kA_k=a'I$ --- solvable to $10^{-16}$ --- but with $a'>a$ strictly:
$a'=3.860$ against $a=3$ at $m=4$, and $4.385$ against $3$ and $5.120$ against $4$ at $m=6$
(\texttt{code/rl\_conj.py}). So for weighted families $F_A(x)=\mu(x+a')$ with $a'\ne a$, and the
identity as stated holds exactly when the two constants agree, which is the projection case
$A_k=P_k$, where $e_2\equiv1$ gives $\sum_kA_k=\operatorname{Adj}(A)=aI$. Everything below is
derived from the wedge coefficients $M_r$ and from $\operatorname{Adj}(A)$ directly and is
unaffected; what does not transfer automatically to the weighted class is a result stated for
mixed characteristic polynomials, since the two constants there are different.

Consequently the band is stated below on the roots of $F_A$, as $|x|\le2\sqrt a$ with $a$ the
constant of $\operatorname{Adj}(A)\preceq aI$, and that is the form the unconditional bounds
establish, through $\lambda^2\le M_1\le am/2$. Written in the $\mu$ variable it becomes
$[a'-2\sqrt a,\,a'+2\sqrt a]$: the radius is set by the $\operatorname{Adj}$ constant and the
centre by the tight one. Both the shift $y=x+a'$ and the bound $|x|\le2\sqrt a$ are checked
directly (\texttt{code/rl\_conj.py}); the two coincide precisely for projection families, where
$a'=a$, and every $\mu$-variable statement of the band in this paper should be read with that
proviso.

The band being sought is the sharp weighted Heilmann--Lieb constant for plane families, and that
constant is not $2\sqrt a$. On projections the conjectural sharp value is $2\sqrt{a-1}$, the
spectral radius of the $a$-regular tree, and $2\sqrt{a-1}<2\sqrt a$ strictly
(\texttt{XuBound.xu\_lt\_note\_target}); it is the case $b=2$ of the bound
$\bigl(\sqrt{a-1}+\sqrt{b-1}\bigr)^2$ on $\operatorname{maxroot}\mu$ conjectured by
Xu\footnote{Personal communication, August 2026, deriving it from a generalisation of
\cite[Conjecture 1]{RL}.}, which is exactly the upper edge of the $(a,b)$-biregular tree band.
So $2\sqrt a$ is the ceiling of the cavity argument used here, not the truth, and an argument of
this kind cannot be sharp. That the sharp value is $(\sqrt{a-1}+\sqrt{b-1})^2$ and not something
smaller is no longer conjectural on the coordinate locus: it is
Proposition~\ref{prop:xusharp}, at every $b$.

In that language the vertex step has an exact form. Setting
$\Theta^{(r)}:=\sum_{|T|=r}M_{\omega_T}\succeq0$ and
$W(x):=\sum_{r\ge1}(-1)^r\Theta^{(r)}x^{\,m-2r}$, one has
$\langle e,W(x)e\rangle=F_A(x)-x\,F_{A^{(e)}}(x)$ for every unit vector $e$, where
$F_A(x)=\mu(x+a)$ and $A^{(e)}$ is the compression to $e^\perp$; consequently
$\operatorname{tr}W=mF_A-xF_A'$ and $F_A'=\sum_iF_{A^{(e_i)}}$ for any orthonormal
basis. The cavity hypothesis is then equivalent to a single matrix inequality: when
$\operatorname{Adj}(A)=aI$, the band $[a-2\sqrt a,\,a+2\sqrt a]$ holds if and only if
\[
W(2\sqrt a)+F_A(2\sqrt a)\,I\ \succeq\ 0 ,
\]
a coordinate-free statement with no recursion in it.

The matrix $W$ has a closed form. If $\omega_T=c_1\wedge\cdots\wedge c_{2r}$ is
decomposable with factor matrix $C=[c_1|\cdots|c_{2r}]$ and Gram matrix $G=C^{\mathsf T}C$,
then $\langle\omega_{\hat\jmath},\omega_{\hat l}\rangle$ is the $(l,j)$ cofactor of
$G$, so
\begin{equation}\label{eq:Momega}
M_{\omega_T}=C\,\operatorname{adj}(G)\,C^{\mathsf T},
\end{equation}
whence $\operatorname{tr}M_{\omega_T}=2r\|\omega_T\|^2$ and
$\operatorname{tr}\Theta^{(r)}=2rM_r$, which is the identity
$\operatorname{tr}W=mF_A-xF_A'$ again. Two of the levels are forced to be scalar:
$\Theta^{(1)}=\operatorname{Adj}(A)=aI$ by hypothesis, and when $m$ is even
$\Theta^{(m/2)}=M_{m/2}\,I$, since a top-degree form is a multiple of the volume
form and $\|\iota_v\omega\|^2=\|\omega\|^2\|v\|^2$ for such. Hence for $m\le4$ every
level is scalar, $W(x)$ is a scalar matrix, and the matrix inequality degenerates to
a scalar one; from $m=5$ on it does not, and that is where the content begins.

Behind the matrix inequality there is a vertex recursion valid in every direction,
which the graph case sees only along the coordinate axes.

\begin{theorem}\label{thm:vertexrec}
Let $A$ be a weighted $2$-plane family on $\mathbb R^m$ with bivectors $\omega_k$,
and let $e$ be a unit vector. Put $f_k=\iota_e\omega_k$, so that
\[
\omega_k=e\wedge f_k+\omega_k',\qquad
f_k,\ \omega_k'\in e^\perp ,
\]
and $\omega_k'$ is the bivector of the compressed family $A^{(e)}$. Then for every
$r\ge1$
\[
M_r(A)=M_r\bigl(A^{(e)}\bigr)
+\sum_{|T|=r}\Bigl\|\sum_{k\in T}f_k\wedge\omega'_{T\setminus k}\Bigr\|^2 ,
\]
and consequently, with $N_r(e)$ denoting that sum of squares,
\[
F_A(x)=x\,F_{A^{(e)}}(x)-N_e(x),\qquad
N_e(x)=\sum_{r\ge1}(-1)^{r-1}N_r(e)\,x^{\,m-2r}.
\]
\end{theorem}

\begin{proof}
Since $e$ is a unit vector, $\iota_e(\omega_k-e\wedge f_k)=f_k-f_k=0$, so
$\omega_k-e\wedge f_k$ lies in $\Lambda^2(e^\perp)$; writing $b_i=\beta_ie+b_i'$
identifies it as $b_1'\wedge b_2'$, the bivector of the compressed block. As
$e\wedge f_k$ is a $2$-form it commutes with the other factors, and $e\wedge e=0$
kills every term using two of them, so
$\omega_T=\omega_T'+e\wedge\sum_{k\in T}f_k\wedge\omega'_{T\setminus k}$. The two
summands lie in $\Lambda^{2r}(e^\perp)$ and $e\wedge\Lambda^{2r-1}(e^\perp)$, which
are orthogonal, so the norms add; summing over $|T|=r$ gives the first display, and
the second is that display multiplied by $(-1)^rx^{\,m-2r}$ and summed.
\end{proof}

We have looked for Theorem~\ref{thm:vertexrec} in the literature and not found it. The
whole arXiv corpus mentioning ``mixed characteristic polynomial'' is seven papers, all
working through the barrier function or interlacing families; the generalizations of
Heilmann--Lieb we can find are to hypergraphs \cite{Zhang}, to spanning subgraphs with
degree constraints, and to multivariate independence polynomials, all of them
graph-theoretic; and a search of arXiv for ``matching polynomial'' together with any of
Grassmann, Pl\"ucker, bivector or exterior algebra returns nothing. That is weak evidence:
the Gram form below is elementary enough to be folklore, the pre-arXiv matching-polynomial
literature is not covered by such a search, and the statement could well appear under
different language, hyperbolic or Lorentzian polynomials, or the theory of stable
polynomials, where Pl\"ucker coordinates do occur but in the different role of
characterizing stability by total nonnegativity. We claim only that we have not found it.

The exterior algebra can be dispensed with entirely, and doing so identifies the
theorem with a known lemma. Let $C$ be the $m\times2r$ matrix whose columns are the
factors of $\omega_T$, so that $\|\omega_T\|^2=\det(C^{\mathsf T}C)$. Splitting off
$e$ is the decomposition $C=eg^{\mathsf T}+C'$ with $g=C^{\mathsf T}e$ and
$C'=(I-ee^{\mathsf T})C$ the compressed factor matrix; since $e^{\mathsf T}C'=0$ and
$\|e\|=1$,
\[
C^{\mathsf T}C=C'^{\mathsf T}C'+gg^{\mathsf T} .
\]
So the content of Theorem~\ref{thm:vertexrec} is the \emph{matrix determinant lemma}
\begin{equation}\label{eq:mdl}
\det\bigl(A+gg^{\mathsf T}\bigr)=\det A+g^{\mathsf T}\operatorname{adj}(A)\,g,
\qquad A=C'^{\mathsf T}C' ,
\end{equation}
and the nonnegativity of the cavity term is nothing but positivity of the adjugate of
a positive semidefinite matrix. In particular $M_r(A)\ge M_r(A^{(e)})$ needs no
computation at all. Splitting $g$ into its $r$ blocks of two identifies the pieces:
the block-diagonal part of the quadratic form \eqref{eq:mdl} is
$\sum_k\theta_kM''_{k,r-1}$ and the off-block-diagonal part is $C_r$, which is how the
cross terms should be read.

For a coordinate family at a coordinate direction $e=e_v$ one has $f_k=\pm e_u$ for
the edges $k=uv$ and $f_k=0$ otherwise, the cross terms vanish, and
$N_r(e_v)=\sum_{u\sim v}m_{r-1}(G-v-u)$: the theorem is then the Heilmann--Lieb
vertex recursion. At a generic direction, or for a family that is not diagonal, the
cross terms do not vanish and the recursion is not the classical one. Two things
follow at once. First, taking $r$ arbitrary, $M_r(A)\ge M_r(A^{(e)})$ for every unit
$e$: the matching numbers of a plane family decrease under compression in every
direction, which for graphs is $m_r(G)\ge m_r(G-v)$. Second, the whole band reduces
to the sign of the cross terms.

\begin{proposition}\label{prop:cross}
Write $X_e:=N_e-\sum_k\theta_kF_{A''_k}$ for the cross-term part of $N_e$, where
$\theta_k=\|f_k\|^2$ and $A''_k$ is the compression to $\{e,f_k\}^\perp$. If
$X_e(x)\le0$ for every weighted $2$-plane family with $\operatorname{Adj}\preceq aI$,
every unit $e$ and every $x\ge2\sqrt a$, then every such family satisfies the band
$[a-2\sqrt a,\,a+2\sqrt a]$.
\end{proposition}

\begin{proof}
Induct on $m$. Put $R_e=F_A/F_{A^{(e)}}$ and $R_k'=F_{A^{(e)}}/F_{A''_k}$. Dividing
Theorem~\ref{thm:vertexrec} by $F_{A^{(e)}}$,
\[
R_e=x-\sum_k\frac{\theta_k}{R_k'}-\frac{X_e}{F_{A^{(e)}}} .
\]
The class is closed under compression, so the inductive hypothesis $R_k'\ge x/2$
applies, and $\sum_k\theta_k=\sum_k\|\iota_e\omega_k\|^2
=\langle e,\operatorname{Adj}(A)e\rangle\le a$. Hence $R_e\ge x-2a/x$, and
$x-2a/x\ge x/2$ exactly when $x\ge2\sqrt a$. The base case $m\le1$ is $R_e=x$.
So $R_e\ge x/2>0$ for $x\ge2\sqrt a$, and $F_A=R_eF_{A^{(e)}}$ has no root there;
$F_A$ is even or odd, so all roots lie in $[-2\sqrt a,2\sqrt a]$.
\end{proof}

The leading cross term can be evaluated exactly, and it is a square.

\begin{theorem}\label{thm:C2}
Let $A$ be a weighted $2$-plane family on $\mathbb R^m$ and $e$ a unit vector, and put
$f_k=\iota_e\omega_k$, $\omega_k'=\omega_k-e\wedge f_k$ and $u_k=\iota_{f_k}\omega_k'$.
Then
\[
C_2=\Bigl\|\sum_kf_k\otimes\omega_k'\Bigr\|^2-\Bigl\|\sum_ku_k\Bigr\|^2 ,
\]
and $\sum_ku_k$ is the off-diagonal $(e,e^\perp)$ block of $\operatorname{Adj}(A)$. So if
$\operatorname{Adj}(A)=aI$ then
\[
C_2=\Bigl\|\sum_kf_k\otimes\omega_k'\Bigr\|^2\ \ge\ 0,
\]
with equality if and only if $\sum_kf_k\otimes\omega_k'=0$.
\end{theorem}

\begin{proof}
For a vector $u$ and forms $\alpha,\gamma$ one has
$\iota_w(u\wedge\alpha)=\langle w,u\rangle\alpha-u\wedge\iota_w\alpha$, hence
\begin{equation}\label{eq:wedgeip}
\langle u\wedge\alpha,\ w\wedge\gamma\rangle
=\langle u,w\rangle\langle\alpha,\gamma\rangle-\langle\iota_w\alpha,\iota_u\gamma\rangle .
\end{equation}
Applying \eqref{eq:wedgeip} to each cross term with $\alpha=\omega_l'$,
$\gamma=\omega_k'$ gives
$C_2=\sum_{k\ne l}\bigl[\langle f_k,f_l\rangle\langle\omega_k',\omega_l'\rangle
-\langle u_k,u_l\rangle\bigr]$.

Now $f_k=\iota_e\omega_k$ lies in the plane of $\omega_k$, and it is orthogonal to $e$,
so it lies in the image of that plane under the projection to $e^\perp$, which is the
plane of $\omega_k'$. Hence $f_k\wedge\omega_k'=0$, and since
$\|f_k\|^2\|\omega_k'\|^2=\|\iota_{f_k}\omega_k'\|^2+\|f_k\wedge\omega_k'\|^2$ for simple
$\omega_k'$, this says exactly
$\|f_k\|^2\|\omega_k'\|^2=\|u_k\|^2$. Therefore
\[
\sum_{k\ne l}\langle f_k,f_l\rangle\langle\omega_k',\omega_l'\rangle
=\Bigl\|\sum_kf_k\otimes\omega_k'\Bigr\|^2-\sum_k\|u_k\|^2,
\qquad
\sum_{k\ne l}\langle u_k,u_l\rangle=\Bigl\|\sum_ku_k\Bigr\|^2-\sum_k\|u_k\|^2,
\]
and subtracting gives the display. For the last claim, polarizing
$\langle v,\Theta_kv\rangle=\|\iota_v\omega_k\|^2$ at $v=e$ against $v\in e^\perp$ gives
$\langle e,\operatorname{Adj}(A)v\rangle=-\langle v,\sum_k u_k\rangle$, so
$\operatorname{Adj}(A)=aI$ forces $\sum_ku_k=0$.
\end{proof}

Theorem~\ref{thm:C2} is sharp in the way the data show: over coordinate families
$\min_eC_2=0$ to machine precision, attained exactly at the coordinate directions, where
each summand $f_k\otimes\omega_k'$ vanishes separately, every $\omega_k$ either
contains $e$ in its plane, so $\omega_k'=0$, or is orthogonal to it, so $f_k=0$. That is
precisely why the classical recursion has no cross terms. Off that locus the minimum is
strictly positive.

The proof reindexes, and doing so gives every level at once.

\begin{theorem}\label{thm:Cr}
For every $r\ge2$,
\[
C_r=\sum_{|S|=r-2}\Bigl[\,Q_S-\|W_S\|^2\,\Bigr],
\qquad
\begin{aligned}
Q_S&=\sum_{k,l\notin S}\langle f_k,f_l\rangle
\bigl\langle\omega_k'\wedge\omega_S',\ \omega_l'\wedge\omega_S'\bigr\rangle,\\
W_S&=\sum_{k\notin S}\iota_{f_k}\bigl(\omega_k'\wedge\omega_S'\bigr),
\end{aligned}
\]
and $Q_S=\bigl\|\sum_{k\notin S}f_k\otimes(\omega_k'\wedge\omega_S')\bigr\|^2\ge0$.
\end{theorem}

\begin{proof}
Writing $T=S\cup\{k,l\}$ reindexes $C_r$ as
$\sum_{|S|=r-2}\sum_{k\ne l\notin S}\langle f_k\wedge\omega_l'\wedge\omega_S',\
f_l\wedge\omega_k'\wedge\omega_S'\rangle$. Apply \eqref{eq:wedgeip} with
$\alpha=\omega_l'\wedge\omega_S'$ and $\gamma=\omega_k'\wedge\omega_S'$. The wedge
$\omega_k'\wedge\omega_S'$ is a product of vectors, hence simple, and
$f_k\wedge(\omega_k'\wedge\omega_S')=(f_k\wedge\omega_k')\wedge\omega_S'=0$ by the proof
of Theorem~\ref{thm:C2}; so
$\|f_k\|^2\|\omega_k'\wedge\omega_S'\|^2=\|\iota_{f_k}(\omega_k'\wedge\omega_S')\|^2$,
and the diagonal terms of the two double sums cancel exactly as before. The Gram matrix
of the vectors $\omega_k'\wedge\omega_S'$ is positive semidefinite, so $Q_S$ is the
squared norm of $\sum_{k\notin S}f_k\otimes(\omega_k'\wedge\omega_S')$.
\end{proof}

At $r=2$ the set $S$ is empty, $W_\emptyset=\sum_k\iota_{f_k}\omega_k'$ vanishes by
tightness, and Theorem~\ref{thm:Cr} is Theorem~\ref{thm:C2}. At $r=3$ it does not
vanish, but it simplifies: since $\iota_{f_n}\omega_n'$ lies in the plane of $\omega_n'$,
the self term $(\iota_{f_n}\omega_n')\wedge\omega_n'$ is zero, and using
$\sum_ku_k=0$,
\[
W_{\{n\}}=\sum_{k\ne n}\omega_k'\wedge\iota_{f_k}\omega_n' .
\]
So the obstruction at $r=3$ is exactly the failure of the compressed planes to be
mutually orthogonal, and it is the only obstruction: everything else is a square.

It is worth being exact about what this settles. Writing the condition $X_e\le0$ at
$x=2\sqrt a$ in the form
\begin{equation}\label{eq:Chat}
\sum_{r\ge2}(-1)^rC_r\,(4a)^{2-r}\ \ge\ 0 ,
\end{equation}
Theorem~\ref{thm:C2} is the leading coefficient of \eqref{eq:Chat}, that is the case
$a\to\infty$. It gives $X_e\le0$ outright when $m\le5$, where $C_2$ is the only cross
term. That range does not improve on the two-moment bounds above, so
Theorem~\ref{thm:C2} does not by itself enlarge the unconditional range; what it supplies
is the first structural reason for the sign.

The obvious next step, bounding $C_3$ by $4a\,C_2$ so as to reach $m\le7$, is not
available, and it fails for a reason worth recording rather than for lack of effort. The
ratio $4a\,C_2/C_3$ decreases steadily in $m$, over cubic graphs it runs $3.49$ at
$m=8$, $2.00$ at $m=10$ (Petersen), $1.38$ at $m=12$ (Franklin), $1.03$ at $m=14$
(Heawood), and crosses below $1$ at $m=16$, where the M\"obius--Kantor graph gives
$0.827$, then $0.687$ (Pappus, $m=18$), $0.590$ (Desargues, $m=20$) and $0.457$ (Nauru,
$m=24$). So $C_3\le4aC_2$ is false in general, and increasingly so.

The reason is visible one level down. For odd $r$ the criterion wants $C_r$ small, that
is $\|W_S\|^2$ close to $Q_S$; but the ratio $\sum_n\|W_n\|^2/\sum_nQ_n$ decays, running
$0.207$, $0.190$, $0.154$, $0.119$, $0.101$ over the same graphs at $m=8,\dots,16$. So
$C_3$ is essentially $\sum_nQ_n$, an undamped sum of Schur squares, and the contraction
$W$ supplies almost no cancellation. Whatever makes \eqref{eq:Chat} hold is therefore the
alternation \emph{between} levels, not any cancellation \emph{within} one.

What the same computation shows is that the truncation is what fails, not the conclusion.
Writing $\Sigma_R$ for the partial sum of \eqref{eq:Chat} over $2\le r\le R$:
\[
\begin{array}{lcrrrrr}
\hline
G & m & \Sigma_2 & \Sigma_3 & \Sigma_4 & \Sigma_5 & \Sigma_6\\
\hline
\text{cube} & 8 & 1.268 & 0.905 & 0.926 & & \\
\text{Petersen} & 10 & 1.624 & 0.812 & 0.891 & 0.891 & \\
\text{Franklin} & 12 & 1.713 & 0.468 & 0.730 & 0.713 & 0.713\\
\text{Heawood} & 14 & 2.007 & 0.060 & 0.679 & 0.603 & 0.607\\
\text{M\"obius--Kantor} & 16 & 2.139 & -0.448 & 0.669 & 0.458 & 0.476\\
\text{Pappus} & 18 & 2.142 & -0.974 & 0.739 & 0.291 & \\
\text{Desargues} & 20 & 2.235 & -1.552 & 0.969 & 0.127 & \\
\text{Nauru} & 24 & 2.231 & -2.656 & 1.786 & & \\
\hline
\end{array}
\]
The values are for one direction per graph, fixed by the seed in
\texttt{code/hl\_Cr.py}, which also regenerates Theorem~\ref{thm:Cr} against the
definition of $C_r$. The terms $|C_r|(4a)^{2-r}$ peak at $r=3$ and then decay
geometrically, at $m=16$ they run
$2.14,\,2.59,\,1.12,\,0.211,\,0.0176$, so the series converges, but the number of
terms needed grows with $m$: three nearly suffice at $m=14$, four at $m=16$, and at
$m=20$ and $24$ the partial sums have not settled by $\Sigma_5$. So \eqref{eq:Chat} is
not a statement of bounded length, and in particular the two-term reading fails outright
from $m=16$. By Theorem~\ref{thm:Cr} every term is a Schur square minus the square of a
contraction, so what a proof must supply is control of the alternating sum of the $Q_S$
across levels.

How much room the inequality actually has is worth measuring, and the answer is that
there is a great deal, and that the recursion cannot use any of it. Evaluating the
exact ratio $R_e=F_A/F_{A^{(e)}}$ at $x=2\sqrt a$ on regular graphs, in units of the
threshold $x/2=\sqrt a$:
\[
\begin{array}{lccc}
\hline
a & \text{graphs} & R_e/(x/2) & \text{spread}\\
\hline
3 & \text{cube, Petersen, Franklin, Heawood} & 1.3662\text{--}1.3692 & 3.0\cdot10^{-3}\\
4 & C_8(1,2),\,C_{10}(1,2),\,C_{12}(1,2) & 1.3488\text{--}1.3497 & 8.6\cdot10^{-4}\\
5 & C_8(1,2,4),\,C_{10}(1,2,5) & 1.3246\text{--}1.3322 & 7.5\cdot10^{-3}\\
\hline
\end{array}
\]
The ratio is a function of $a$ alone to three digits, independent of the graph and of the
dimension, and exceeds the threshold by a third. The slack $-X_e/(F_{A^{(e)}}\sqrt a)$
likewise \emph{increases} with $m$, running $0.033,0.044,0.049,0.058$ over the four cubic
graphs above.

That constant can be identified exactly. Averaging the vertex identity over directions
gives $\mathbb E_eF_{A^{(e)}}=F_A'/m$, which is $F_A'=\sum_iF_{A^{(e_i)}}$ again, so
the direction-averaged ratio is
\[
\bar R=\frac{mF_A}{F_A'}=\Bigl(\frac1m\sum_i\frac1{x-\rho_i}\Bigr)^{-1}=\frac1{G(x)},
\]
$G$ the Stieltjes transform of the root measure of $F_A$. For a coordinate family $F_A$ is
the matching polynomial, and the Kesten--McKay measure of parameter $a$ has Stieltjes
transform $G(z)=2(a-1)/\bigl((a-2)z+a\sqrt{z^2-4(a-1)}\bigr)$. At $z=2\sqrt a$ the radical
is \emph{exactly} $2$, and everything collapses:

\begin{proposition}\label{prop:km}
$\dfrac{1}{G(2\sqrt a)}\Big/\sqrt a=\dfrac{\sqrt a+2}{\sqrt a+1}=1+\dfrac1{1+\sqrt a}$.
\end{proposition}

The identity is proved below by the finite and effective route described next; it is not
formalized.

The passage from a finite graph to that value is usually made by citing convergence of the
matching measure. It need not be: the evaluation point lies strictly outside the support,
so the whole step can be made finite and effective, with an explicit rate. We are grateful
to P\'eter Csikv\'ari for pointing us to the mechanism, which is his with P.\ E.\ Frenkel
\cite{CsF}, after Ab\'ert and Hubai for the chromatic polynomial.

Their theorem is that for a monic multiplicative graph polynomial of bounded exponential
type, the power sum $p_j(G)$ of the roots is a fixed linear combination of counts of
\emph{connected} subgraphs of $G$; the matching polynomial is such a polynomial. In
Godsil's formulation \cite{Go} the matching-measure moments count closed \emph{tree-like}
walks, and the trace of such a walk of length $j$ carries at most $j/2$ edges. Since a
subgraph with fewer than $\operatorname{girth}(G)$ edges is a forest, and since in an
$a$-regular graph the number of copies of a fixed tree is $|V(G)|$ times a function of $a$
alone, one gets an identity rather than a limit.

\begin{proposition}\label{prop:km-moments}
Let $G$ be $a$-regular of girth $g$. Then $p_j(G)/|V(G)|$ equals the $j$-th moment of the
Kesten--McKay measure of parameter $a$ for every $j<2g$. The threshold is sharp: for
$K_4$, $K_5$, $K_{3,3}$, $K_{4,4}$, the cube, Petersen and Heawood, the first differing
moment is exactly $2g$ whenever one exists at all
\textup{(}\texttt{code/moment\_girth.py}, exact arithmetic\textup{)}.
\end{proposition}

The agreement of moments below $2g$ is proved, by the tree-like walk count of \cite{Go}. The
sharpness of the threshold is measured on the seven graphs named, not proved.

Because Heilmann--Lieb confines the roots to $[-\rho,\rho]$ with $\rho=2\sqrt{a-1}$, while
the evaluation point is $x=2\sqrt a$, the ratio $\rho/x=\sqrt{1-1/a}$ is strictly below
one. A finite geometric expansion with remainder, and no infinite series anywhere, then
converts matching moments into matching transforms.

\begin{proposition}\label{prop:transfer}
Let $\nu_1,\nu_2$ be probability measures on $[-\rho,\rho]$ whose moments agree up to order
$k$, and let $x>\rho$. Then
\[
\bigl|G_{\nu_1}(x)-G_{\nu_2}(x)\bigr|\ \le\ \frac{2\rho^{k+1}}{x^{k+1}(x-\rho)},
\]
and if both transforms are at least $c>0$ the reciprocals differ by at most that quantity
divided by $c^2$.
\end{proposition}

\begin{proof}
For $|t|\le\rho<x$ one has the identity
$1/(x-t)=\sum_{j\le k}t^j/x^{j+1}+t^{k+1}/\bigl(x^{k+1}(x-t)\bigr)$, whose right-hand side
does not depend on $k$. Integrating, the polynomial parts agree by hypothesis, and each
remainder is at most $\rho^{k+1}/\bigl(x^{k+1}(x-\rho)\bigr)$ in absolute value. The
statement for reciprocals is $|1/S_1-1/S_2|=|S_2-S_1|/(S_1S_2)$.
\end{proof}

\begin{remark}\label{rem:csik}
Csikv\'ari asks (personal communication, August 2026) whether $\mu_G(2\sqrt{d-1})$ has a
combinatorial meaning for $d$-regular $G$ with $|V(G)|$ even, observing that it should at
least be a positive integer. It does, and the answer is in his own paper with Bencs
\cite{BCs}. Their polynomial
$R_G(z)=\sum_{M}(-z)^{|M|}\prod_{v\notin V(M)}(z+d_v-1)$ is, for $d$-regular $G$, a
transformation of the matching polynomial: comparing coefficients termwise,
\[
R_G(z)=z^{n/2}\,\mu_G\!\Bigl(\frac{z+d-1}{\sqrt z}\Bigr).
\]
The point $2\sqrt{d-1}$ is not an arbitrary argument of that substitution. Solving
$(\sqrt z+ (d-1)/\sqrt z)=2\sqrt{d-1}$ gives $(\sqrt z-\sqrt{d-1})^2=0$, so it is exactly
the \emph{double root}, attained only at $z=d-1$. Their Corollary 2.7 expands
$R_G(w)=\sum_{A\in\mathcal{PF}}2^{c(A)}(w-1)^{n-|A|}$ over pseudo-forests $A$, edge sets each
of whose components carries at most one cycle, with $c(A)$ the number of cycles. Putting
$w=d-1$,
\begin{equation}\label{eq:pf}
\mu_G\bigl(2\sqrt{d-1}\bigr)\;=\;(d-1)^{-n/2}\sum_{A\in\mathcal{PF}(G)}2^{c(A)}(d-2)^{\,n-|A|}.
\end{equation}
At $d=3$ the weight is $1$ and \eqref{eq:pf} reads: $2^{n/2}\mu_G(2\sqrt2)$ is the number of
pseudo-forests of $G$ counted with multiplicity $2$ per cycle. The values are $19$, $74$,
$76$, $297$, $1162$ for $K_4$, $K_{3,3}$, the prism, the cube and Petersen.

The transformation, the double root and the evaluation are machine-checked in
\texttt{RamaLean/PseudoForest.lean}; Corollary 2.7 is \cite{BCs} and enters as a hypothesis.
Identity \eqref{eq:pf} is verified by brute force in \texttt{code/pseudoforest.py} on nine
regular graphs of degrees $3$ and $4$ and orders $4$ to $10$, in exact arithmetic. The
positivity and integrality Csikv\'ari notes in passing are separately machine-checked in
\texttt{RamaLean/EvenEval.lean}, together with the observation that the parity hypothesis is
needed: for $|V(G)|$ odd the value is $2\sqrt{d-1}$ times an integer.
\end{remark}

Propositions~\ref{prop:km-moments} and~\ref{prop:transfer} together bound the distance from
the finite-graph ratio to the value of Proposition~\ref{prop:km} by a quantity decaying
like $(1-1/a)^{g}$, with no convergence theorem invoked. The bound is not tight, since
$|m_j|\le\rho^j$ discards the $j^{-3/2}$ decay of the Kesten--McKay moments at the edge: at
$a=3$ it gives $0.62$, $0.41$, $0.28$ at girths $4,5,6$ against observed errors of
$3.2\cdot10^{-3}$, $9.2\cdot10^{-4}$, $2.2\cdot10^{-4}$. What it does give, and the
citation did not, is an effective statement about a finite graph.

So the cavity ratio exceeds the threshold the induction can propagate by exactly
$1/(1+\sqrt a)$, and the excess tends to $0$ as $a$ grows, which is the analytic shadow
of $2\sqrt a$ being asymptotically attained.

The limit appears to be attained from above, never crossed, which suggests that
Proposition~\ref{prop:km} states an infimum rather than merely a limit:

\begin{conjecture}\label{conj:km}
For every weighted $2$-plane family with $\operatorname{Adj}(A)=aI$ and every unit $e$,
\[
\frac{R_e}{x/2}\ \ge\ 1+\frac1{1+\sqrt a}\qquad\text{at }x=2\sqrt a .
\]
\end{conjecture}

This was not violated in any family tested, over graphs and over general weighted plane
families, minimizing over directions; the smallest margin observed was
$2.2\cdot10^{-4}$, at the highest-girth cubic graph, which is where
Proposition~\ref{prop:km} says it should be. Conjecture~\ref{conj:km} is strictly
stronger than $X_e\le0$, which is its $c=1$ case, and by Proposition~\ref{prop:rigid} it
cannot be proved by this induction either, the recursion propagates $c=1$ and nothing
else, whatever is true. Against the exact values the prediction is
$1.366025$ at $a=3$, $1.333333$ at $a=4$ and $1.309017$ at $a=5$, and the error decays
with the girth exactly as Kesten--McKay convergence requires: over cubic graphs
$3.2\cdot10^{-3}$ at girth $4$ (the cube), $9.2\cdot10^{-4}$ at girth $5$ (Petersen), and
$2.2\cdot10^{-4}$, $1.7\cdot10^{-4}$, $1.2\cdot10^{-4}$, $1.2\cdot10^{-4}$ at girth $6$
(Heawood, M\"obius--Kantor, Pappus, Desargues). None of that slack is available to the
argument:

\begin{proposition}\label{prop:rigid2}
Fix $a>0$ and suppose the inductive hypothesis of Proposition~\ref{prop:cross} is
strengthened from $R\ge x/2$ to $R\ge c\,(x/2)$ for some $c>0$. At $x=2\sqrt a$ the
recursion sustains this only for $c=1$.
\end{proposition}

\begin{proof}
Propagating $R\ge c\sqrt a$ requires $c\sqrt a\le2\sqrt a-a/(c\sqrt a)$, which after
clearing the denominator is $a(c-1)^2\le0$, so $c=1$.
\end{proof}

There is one hypothesis the argument has not used, and it changes what has to be proved.
The class is \emph{not} closed under compression in the tight sense:

\begin{proposition}\label{prop:deficient}
For a weighted $2$-plane family with $\operatorname{Adj}(A)=aI$ and any unit $e$,
\[
\operatorname{Adj}\bigl(A^{(e)}\bigr)=a\,I_{e^\perp}-F,\qquad F=\sum_kf_kf_k^{\mathsf T},
\qquad \operatorname{tr}F=a .
\]
\end{proposition}

\begin{proof}
For $u,v\perp e$ one has $\iota_u\omega_k=-\langle u,f_k\rangle e+\iota_u\omega_k'$ with
$\iota_u\omega_k'\perp e$, so
$\langle\iota_u\omega_k,\iota_v\omega_k\rangle
=\langle u,f_k\rangle\langle v,f_k\rangle+\langle\iota_u\omega_k',\iota_v\omega_k'\rangle$;
summing and using $\operatorname{Adj}(A)=aI$ gives the display, and
$\operatorname{tr}F=\sum_k\theta_k=\langle e,\operatorname{Adj}(A)e\rangle=a$.
\end{proof}

So compressing costs exactly a positive semidefinite matrix of trace $a$, and the children
of a tight family are deficient by a definite amount. Spending that deficiency strengthens
the induction. Write $d(e)=a-\langle e,\operatorname{Adj}(A)e\rangle\ge0$ for the local
deficiency, and replace the hypothesis $R\ge x/2$ of Proposition~\ref{prop:cross} by
\[
(\mathrm H)\qquad R_e\ \ge\ \frac x2+\frac{2\,d(e)}{x}.
\]
This is consistent at the base ($m\le1$: $R_e=x$, $d=a$, and $x/2+2a/x=x$ exactly), and at
a tight family $d\equiv0$, so it still delivers the band.

\begin{proposition}\label{prop:defstep}
Let $g_k=\langle\hat f_k,F\hat f_k\rangle$ be the deficiency the children inherit. Then
$(\mathrm H)$ propagates provided
\[
X_e\ \le\ \frac{2F_{A^{(e)}}}{x}\sum_k\frac{\theta_k\,g_k}{a+g_k},
\]
and the estimate is lossless: the two sides of the resulting comparison are identically
equal.
\end{proposition}

Proposition~\ref{prop:defstep} is proved below and is not formalized; the margins quoted for it
afterwards are measured.

The right-hand side is \emph{strictly positive}, since $g_k\ge\theta_k>0$. So the
strengthened induction does not require $X_e\le0$; it requires $X_e$ below an explicit
positive quantity. On the graphs tested $(\mathrm H)$ held with margin $0.52$ to $0.71$,
and the sufficient condition held with the right-hand side exceeding $|X_e|$ by a factor
of five to twenty.

That is worth having because an upper bound on $X_e$ does exist, and the deficiency
supplies it. Applying the vertex identity a second time writes each child polynomial as
$F_{A''_k}=\bigl(F_{A^{(e)}}-q_k\bigr)/x$ with
$q_k=\langle\hat f_k,W^{(e)}(x)\hat f_k\rangle$, so
$\sum_k\theta_kF_{A''_k}=\bigl(aF_{A^{(e)}}-\operatorname{tr}(FW^{(e)})\bigr)/x$, and
substituting into $X_e=N_e-\sum_k\theta_kF_{A''_k}$ gives a relation between the vertex
matrices at two consecutive levels:

\begin{proposition}\label{prop:twolevel}
$x\,X_e=\operatorname{tr}\bigl(F\,W^{(e)}(x)\bigr)-x\,\langle e,W(x)e\rangle
-a\,F_{A^{(e)}}(x)$, and consequently
\[
x\,X_e\ \le\ a\,\lambda_{\max}\bigl(W^{(e)}(x)\bigr)-x\,\langle e,W(x)e\rangle
-a\,F_{A^{(e)}}(x).
\]
\end{proposition}

\begin{proof}
The identity is the substitution above. For the bound, $F=\sum_k\theta_k\hat f_k\hat
f_k^{\mathsf T}$ is a nonnegative combination of unit rank-ones with
$\sum_k\theta_k=\operatorname{tr}F=a$, so
$\operatorname{tr}(FW^{(e)})=\sum_k\theta_kq_k\le\bigl(\sum_k\theta_k\bigr)
\lambda_{\max}(W^{(e)})$; no spectral theory is needed beyond the definition of
$\lambda_{\max}$.
\end{proof}

This is the first upper bound on $X_e$ in the development, every earlier statement
bounded it below or computed it exactly, and the trace step is not lossy where it
matters: on $K_5$ and $K_6$, where $X_e=0$, it holds with equality. Together with
Proposition~\ref{prop:defstep}, whose target is strictly positive, the band is therefore
reduced to a spectral bound on the child's vertex matrix.

That bound has a clean form in the cavity ratios, and the form explains what is left.
Write $r=R_e/(x/2)$ and $r'=\bigl(\max_{e'}R'_{e'}\bigr)/(x/2)$. Since
$\lambda_{\max}(W^{(e)})=F_{A^{(e)}}-x\min_{e'}F_{A''}$ and
$\min_{e'}F_{A''}=F_{A^{(e)}}/\max_{e'}R'_{e'}$, the bound
$\lambda_{\max}(W^{(e)})\le xF_A/a-3F_{A^{(e)}}$ is exactly
\begin{equation}\label{eq:ratiotarget}
r'\ \le\ \frac1{2-r}.
\end{equation}

\begin{proposition}\label{prop:kmtarget}
The Kesten--McKay value of Proposition~\ref{prop:km} satisfies \eqref{eq:ratiotarget}
automatically: with $u=1/(1+\sqrt a)$ and $r=r'=1+u$, one has $(1+u)(1-u)=1-u^2\le1$, and
the slack is exactly
\[
\frac1{1-u}-(1+u)=\frac{u^2}{1-u}=\frac1{\sqrt a\,(1+\sqrt a)} .
\]
\end{proposition}

So the target is not merely consistent with the Kesten--McKay picture, it is
\emph{implied} by it, with a slack that vanishes like $1/a$, matching the asymptotic
tightness of $2\sqrt a$. On the graphs tested \eqref{eq:ratiotarget} holds throughout, and
it is tight exactly where $X_e=0$: at $K_5$ and $K_6$ the two sides agree to four digits
($1.5673$ and $1.5194$), while the larger-girth graphs have margin. What remains
unproved is therefore not \eqref{eq:ratiotarget} in the abstract but the gap between the
true ratios and their Kesten--McKay values, which is Conjecture~\ref{conj:km} together
with its analogue one level down.

The constant is therefore not an artefact of the estimate: $\sqrt a$ is the fixed
point of $\rho\mapsto x-a/\rho$ at $x=2\sqrt a$, the value at which the two roots of
$\rho^2-x\rho+a$ collide. For $x>2\sqrt a$ the inductive hypothesis has an interval
of admissible thresholds and the argument has slack; at $x=2\sqrt a$ the interval is
a point and it has none. That is exactly why $2\sqrt a$, and not the conjectured
$2\sqrt{a-1}$, is the ceiling of every argument of this shape. The cross terms
$C_r=N_r-\sum_k\theta_kM''_{k,r-1}
=\sum_{|T|=r}\sum_{k\ne l\in T}\langle f_k\wedge\omega'_{T\setminus k},
f_l\wedge\omega'_{T\setminus l}\rangle$ were nonnegative in every family tested, and at
$r=2$ that is now Theorem~\ref{thm:C2}. The nonnegativity is not termwise, and it is
worth recording that it is not, since the termwise statement is the natural thing to try
and it is false: single pairs $(k,l)$
and even whole summands $\sum_{k\ne l\in T}$ take negative values, the smallest
observed being $-1.9\cdot10^{-3}$ and $-3.8\cdot10^{-3}$ over weighted plane
families in dimensions six and eight. If $C_r\ge0$ holds it is therefore a statement about the sum over all
$T$ of a given size, not about its terms. Summed, it survives minimization over the
sphere, and sharply: for coordinate families $\min_e C_2=0$ to machine precision,
attained exactly at the coordinate directions, where the classical recursion lives,
while for noncommuting projection and general plane families the minimum is strictly
positive, growing from $0.048$ at $m=4$ to $0.777$ at $m=8$. Finally
$X_e=\sum_{r\ge2}(-1)^{r-1}C_rx^{\,m-2r}$ was nonpositive at $x=2\sqrt a$ in every case
under direct maximization over the sphere, zero to $4\cdot10^{-12}$ for coordinate
families, strictly negative otherwise, over graph families in dimensions $4,5,6$,
projection families in dimensions $4,6$ and weighted plane families in dimension $4$.
The identity of Theorem~\ref{thm:vertexrec} was checked to $10^{-13}$ on graph,
projection and general weighted plane families at generic directions
(\texttt{code/hl\_Wspec.py}, \texttt{code/hl\_Xsign.py}).

Both of the bounds that follow are rungs of a single ladder, and it is worth saying what the
ladder reaches before giving them. Writing $y=x^2$, the $M_r$ are the elementary symmetric
functions of the $y$-roots, so with $p_k$ the power sums one has $y_{\max}^k\le p_k$, and any
a priori bound $p_k\le(4a)^k$ gives the band (\texttt{MomentLadder.le\_of\_power\_sum\_le}).
Writing $p_k=(m/2)c_k^k$, the $k$-th rung covers dimension $m$ exactly when
\begin{equation}\label{eq:reach}
   m\ \le\ 2\Bigl(\frac{4a}{c_k}\Bigr)^{k}
\end{equation}
(\texttt{reach\_of\_moment}). Measured on the class, $c_k^k$ is the number of closed walks of
length $2k$ from a root of the $a$-regular tree, so $c_1=a$ and $c_2=\sqrt{a(2a-1)}$, which are
the two rungs below; that count grows like $\bigl(4(a-1)\bigr)^kk^{-3/2}$, the base being the
square of the tree's spectral radius, against a target of $(4a)^k$. Since $4(a-1)<4a$, the reach
\eqref{eq:reach} is unbounded in $k$ (\texttt{reach\_unbounded}): at $a=3$ it is $146$, $509$ and
$1489$ at $k=5,7,9$, against the $8$ and $18$ of the first two rungs (\texttt{code/moments2.py}).
What the reach costs is an a priori bound on $p_k$, and the natural candidate is
$p_k\le(m/2)W_{2k}$ with $W_{2k}$ the tree walk count; for graphs that is Godsil's path-tree
argument, the path tree being a subtree of the universal cover. Off the coordinate case it is
false. Weighted families with $\operatorname{Adj}(A)=aI$ exceed it by $1.008$, $1.134$ and
$1.315$ at $m=6,8,10$ and $a=3$, the excess growing in both $k$ and $m$, so no constant-factor
version survives either; coordinate families obey it, with equality at $k=1,2$
(\texttt{code/momentbound.py}, \texttt{code/p16verify.py}). Weakening it to a rate,
$p_k\le(m/2)R^kW_{2k}$, leaves the reach unbounded exactly when $R<a/(a-1)$
(\texttt{rate\_lt\_target}, \texttt{reach\_unbounded\_of\_rate}), and the measured rates sit far
inside that: $0.97$, $1.02$, $1.05$, $1.07$, $1.09$ at $m=6,\dots,14$ and $a=3$
(\texttt{code/excessrate.py}). That is still not a route, for a reason that is structural rather
than numerical.

\begin{proposition}\label{prop:ladder-nogo}
Let $y_1,\dots,y_n\ge0$, let $c>0$, and suppose $\sum_iy_i^k\le Cc^k$ for every $k\ge1$, with
$C$ independent of $k$. Then $y_i\le c$ for every $i$.
\end{proposition}

\begin{proof}
If $y_i>c$ then $(y_i/c)^k\le\bigl(\sum_jy_j^k\bigr)/c^k\le C$ for every $k$, while the left
side is unbounded \textup{(}\texttt{MomentLadder.band\_of\_all\_moments}\textup{)}.
\end{proof}

Applied with $c=4(a-1)R$, a rate bound holding at every rung already forces
$y_{\max}\le4(a-1)R$. So at $R<a/(a-1)$ the hypothesis is strictly stronger than the band it is
meant to prove, and at $R\ge a/(a-1)$ every rung collapses to $m\le2$
(\texttt{reach\_trivial\_of\_rate}). The ladder has content only at a single finite $k$, where it
buys the finite range \eqref{eq:reach} and nothing more; the dimension restriction is intrinsic
to it.

Two bounds follow unconditionally from real-rootedness alone. Writing
$F_A=x^m-M_1x^{m-2}+M_2x^{m-4}-\cdots$, one has $M_1=\tfrac12\operatorname{tr}
\operatorname{Adj}(A)\le am/2$ and, using $\|\omega_j\wedge\omega_l\|^2\ge
1-\operatorname{tr}(P_jP_l)$, $M_1^2-2M_2\le\operatorname{tr}
\operatorname{Adj}(A)^2-\sum_kc_k^2$. Since the largest root $\lambda$ satisfies
$\lambda^2\le M_1$ and $\lambda^4\le M_1^2-2M_2$, the band
$[a-2\sqrt a,a+2\sqrt a]$ holds whenever $m\le8$, and whenever
$m\le16+\bigl(\sum_kc_k^2\bigr)/a^2$, which for projection families is
$m\le32a/(2a-1)$. The second reproduces the known two-moment bound for projections
and extends it to all weighted plane families. Adversarial search over rank-two projection families and
over random biregular graphs reaches $0.83$ of the threshold at $a=3$ and $0.72$ at
$a=4$, with no violation; graphs are marginally more extremal than general
projections.

\subsection{The reflection for projections}

At $b=2$ the whole question also collapses onto the known half, by a symmetry.

\begin{proposition}\label{prop:reflect}
Let $P_1,\dots,P_q$ be rank-$2$ orthogonal projections on $\mathbb R^p$ with
$\sum_kP_k=aI$. Then $\mu[P_1,\dots,P_q](2a-y)=(-1)^p\mu[P_1,\dots,P_q](y)$; in
particular $\lambda_{\min}+\lambda_{\max}=2a$.
\end{proposition}

\begin{proof}
The mixed characteristic polynomial is affine-linear in each slot, so it commutes
with taking expectations slotwise, and for rank-one slots
$\mu[v_1v_1^{\mathsf T},\dots]=\det\bigl(yI-\sum_kv_kv_k^{\mathsf T}\bigr)$, since a
multi-affine polynomial is sent by $\prod_k(1-\partial_k)$ to its value at
$z=(-1,\dots,-1)$. Choose an orthonormal pair $e_k,f_k$ spanning the range of $P_k$
and set $v_k=e_k+\varepsilon_kf_k$ with $\varepsilon_k$ independent uniform signs.
Then $\mathbb E\,v_kv_k^{\mathsf T}=P_k$ and
$v_kv_k^{\mathsf T}=P_k+\varepsilon_kS_k$ with
$S_k=e_kf_k^{\mathsf T}+f_ke_k^{\mathsf T}$ and $S_k^2=P_k$, so
$\mu(y)=\mathbb E_\varepsilon\det\bigl((y-a)I-\sum_k\varepsilon_kS_k\bigr)$. That is
an even or odd function of $t=y-a$ according to the parity of $p$, because
$\varepsilon\mapsto-\varepsilon$ preserves the distribution.
\end{proof}

The last step, which carries the whole proposition, is formalized in
\texttt{RamaLean/SignReflection.lean} as \texttt{sign\_avg\_reflect}: for any family of
matrices $S_1,\dots,S_q$ the sign average
$\sum_{\varepsilon\in\{\pm1\}^q}\det\bigl(tI-\sum_k\varepsilon_kS_k\bigr)$ satisfies
$\rho(-t)=(-1)^n\rho(t)$, with no hypothesis on the $S_k$; the shift and the exchange
of band edges $2a-(\sqrt{a-1}+1)^2=(\sqrt{a-1}-1)^2$ are
\texttt{mixedChar\_reflect} and \texttt{band\_edges\_swap}. The step identifying the
sign average with $\mu$ is formalized in the same file: by \texttt{sum\_prod\_sgn} a
nonempty sign monomial averages to zero over $\{\pm1\}^q$, proved by the same
involution, so by \texttt{sum\_multiaffine} the average of a multi-affine function
returns $2^q$ times its constant term, which for
$\det\bigl(yI+\sum_kz_kA_k\bigr)$ with rank-one $A_k$ is what
$\prod_k(1-\partial_{z_k})$ extracts.

Since $2a-(\sqrt{a-1}+1)^2=(\sqrt{a-1}-1)^2$, the two band edges are exchanged by
the reflection, so at $b=2$ the lower bound is \emph{equivalent} to the upper one.
For diagonal families the upper bound is Godsil's path-tree theorem and $b=2$ is
therefore settled, that is Theorem~\ref{thm:subdiv} again. For projections it is
not: the path tree is a graph construction, and the only bound available in that
generality is the Marchenko--Pastur edge $(\sqrt a+\sqrt b)^2$, which is strictly
larger than $(s+t)^2$. So at $b=2$ the projection statement is reduced to its upper
half, not closed. The mechanism in the diagonal case is transparent: there
$\sum_k\varepsilon_kS_k$ is a signed adjacency matrix, the expectation is
Godsil--Gutman, and one recovers $\nu_{S(H)}(y)=\mu_H(y-a)$.

This also locates the obstruction at $b\ge3$ sharply. Writing
$N_k=v_kv_k^{\mathsf T}-P_k$ one has $N_k^2=(b-2)N_k+(b-1)P_k$, with spectrum
$\{b-1,\,-1^{(b-1)}\}$ on the range of $P_k$; the reflection needs this invariant
under $N_k\mapsto-N_k$, which happens exactly when $b=2$. Equivalently: the mean
root is always $a$, since $\sum_i\lambda_i=\operatorname{tr}\sum_kP_k=pa$, while the
band centre is $a+b-2$. The conjecture therefore asserts a skew of exactly $b-2$,
which is the operator form of the displacement in Theorem~\ref{thm:subcoef} and is
invisible to any reflection or first-moment argument.

The source of that skew can be named exactly. The jump $N_k$ has spectrum
$\{b-1,\,(-1)^{(b-1)}\}$ on the range of $P_k$, whose moments are
\[
m_1=0,\qquad m_2=b(b-1),\qquad m_3=b(b-1)(b-2).
\]
So the third moment vanishes precisely when $b=2$: Proposition~\ref{prop:reflect}
is not an accident of small $b$ but the vanishing of the jump skewness, and for
$b\ge3$ any argument reaching the lower edge must see a third-order quantity.

\section{What this leaves}\label{sec:leaves}

The biregular case is not a leftover from the refutation; it is where the phenomenon actually
lives, and this paper has tried to say how much room it has.

On the positive side the case is genuinely constrained. Complete bipartite graphs never refute
it, by a margin $\sqrt{d-1}-h_d/2$ that Gershgorin on the Hermite Jacobi matrix makes positive
for every $d$, and the ratio route closes the inner edge on a nontrivial class. Those are the two
proved statements here, and they are not vacuous: the first covers a family whose gap can be made
arbitrarily wide, and the second reaches the edge exactly, its certificate degenerating at the gap
edge and nowhere earlier.

On the negative side the margin decays like $n^{-2/3}$. The statement is true but tight, and that
single fact disqualifies most of what one would try next, including the Gershgorin constant that
proves the complete bipartite case: any argument yielding a bound independent of the size is known
in advance to be insufficient. What is needed is a Friedman-type edge theorem, in which the roots
fill out $\operatorname{spec}(T)$ without escaping it, at the soft-edge scale.

\label{sec:obstructions}
The obstructions are the part we would most want a reader to take away, because each one closes a
route that looks open from outside.

\begin{itemize}
\item No bound of the form $\operatorname{maxroot}\mu\le F(a,b,p,q,\operatorname{spec}A_1,\dots,
  \operatorname{spec}A_q)$ can reduce the problem. Every $A_k$ in a tight rank-$b$ projection
  family has spectrum $\{1^{(b)},0^{(p-b)}\}$, so any such $F$ is constant on the class; the class
  contains commuting families whose greatest root tends to the band edge, so $F$ is at best that
  edge, which is the statement to be proved. Both halves are machine-checked in
  \texttt{RamaLean/SpectralNoGo.lean}: that every functional factoring through the power traces of
  the blocks agrees on any two families of idempotents with equal block traces
  (\texttt{const\_of\_factors\_through\_traces}, from $A^{j+1}=A$ and nothing else), and that an $F$
  which is both valid on the class and strong enough to give the target satisfies $F=L$ exactly, so
  that the bound it yields is the target verbatim (\texttt{spectral\_bound\_is\_the\_target}). This
  is why the barrier route yields no reduction, and
  it applies equally to the Poisson-binomial variance of \S\ref{sec:rank}, which is a spectral
  functional and is constant at $(b/a)(1-1/a)$ on the class
  (\texttt{variance\_of\_idem}). A proof must read a joint invariant of
  the family. The commutator is one: over random tight projection families the greatest root
  correlates with $\sum_{i<j}\|[A_i,A_j]\|_F^2$ at $-0.33$, $-0.50$ and $-0.73$ for $C_4$, $C_6$
  and $C_8$, and is largest at zero, which is the commuting locus
  (\texttt{code/jointinv.py}). That correlation has since been tested adversarially rather than on
  random families: hill climbing inside the projection class to \emph{maximise} the greatest root,
  each block rotated by the exponential of a random antisymmetric generator so that idempotency
  stays exact, from $76$ random starts across $C_4$, $C_6$, $K_{3,3}$ and the cube at $600$ steps
  each. It accepted $7661$ uphill steps and never reached the commuting value, falling short by
  $0.000031$, $0.000857$, $0.002763$ and $0.010813$ respectively, with the commutator at the best
  point below its starting value in every family (\texttt{code/jointinv.py}, full configuration;
  \texttt{-{}-quick} truncates it to two families and $120$ steps). It is still not a theorem, and
  the functional it names is the one a proof would have to read.

  The obvious local form of that statement is false, and its failure gives the right one. Both
  quantities are quadratic forms on the kernel at a commuting family: the greatest root has the
  curvature form $\Lambda$, and since $[A_i,A_j]=\varepsilon([P_i,D_j]+[D_i,P_j])+O(\varepsilon^2)$
  the commutator has $\Gamma(D)=\sum_{i<j}\|[P_i,D_j]+[D_i,P_j]\|_F^2$, a sum of squares. A Lyapunov
  statement would read $\Lambda\le-c\,\Gamma$ with $c>0$, and that forces
  $\ker\Lambda\subseteq\ker\Gamma$. It fails: $\ker\Gamma$ is exactly the conjugation orbit, of
  dimension $\binom n2$, while $\ker\Lambda$ is strictly larger at $b=2$, of dimension $8$ against
  $6$ at $C_4$ and $24$ against $15$ at $C_6$, and directions on the cone with $Q=0$, $\Lambda=0$ and
  $\Gamma>0$ are exhibited at both. So no inequality of that shape holds.
  What replaces it is visible along those same directions: the root falls at order four where the
  commutator grows at order two, and the two exponents pair into a constant. At $C_4$ the ratio
  $(y_0-\lambda_{\max})/C^2$ holds at $2.9\times10^{-4}$ to two digits over
  $t\in[0.02,0.20]$ while $(y_0-\lambda_{\max})/t^4$ holds near $9\times10^{-7}$
  (\texttt{code/lyapunov.py}). The shape to aim at is therefore
  $\lambda_{\max}\le y_0-c\,C^2$, quadratic in the commutator rather than linear, and the degenerate
  directions that make the second-order test blind are what distinguishes the two.

  Those directions also carry the sharpest local test the extremality claim has been put to, and it
  passes. On $\ker\Lambda$ the curvature form says nothing, so whether the commuting family is even a
  local maximum there is decided at order four; a direction with $c_4\le0$ would refute the claim
  outright. Over $20$ directions of $\ker\Lambda$ on the cone at each of $C_4$ and $C_6$, found by
  solving $Q=0$ and $\Gamma=1$ inside $\ker\Lambda$ from random starts, the root falls at every
  direction and every step size, with $\ker\Lambda$ of dimension $8$ against an orbit of $6$ at
  $C_4$ and $24$ against $15$ at $C_6$. The resolvable fourth-order constants run from
  $7.7\times10^{-5}$ to $1.73$ at $C_4$ and from $2.8\times10^{-2}$ to $3.16$ at $C_6$; three
  directions at $C_4$ are reported as unresolved rather than quoted, their $y_0-\lambda_{\max}$
  falling to about $4\times10^{-12}$, which is the precision floor of the root solve. A direction
  does not determine a curve, order two leaving the off-diagonal blocks of the correction free, so
  the sign is also checked along curves whose second-order term is perturbed inside the kernel:
  over $240$ curves at each family, twenty directions carrying twelve perturbed corrections apiece,
  the smallest constant found is $1.2\times10^{-6}$ at $C_4$ and $2.3\times10^{-2}$ at $C_6$, both
  positive (\texttt{code/quartic.py}). Since $\Gamma$ is normalised to $1$ on these directions the
  same constants are $c_4/\Gamma^2$, which is the constant in the quadratic shape.
  The shape is local of necessity and not by caution: $y_0-\lambda_{\max}$ is bounded on the class
  while $C^2$ is not, so the ratio falls as the commutator grows and no global constant can work.
\item A barrier argument that sees only the trace and $\sum_kP_k=aI$ is reading the
  Marchenko--Pastur law, and $A_k=(b/p)I_p$ satisfies both while exceeding the tree band
  already at $p=4$. No such argument can reach either edge.
\item Compressions cannot reach the inner end at all: a compression takes values the original
  excludes, so no argument phrased in Rayleigh quotients can exclude them.
\item The coefficient ladder's reach is unbounded in the moment order, but a bound holding at
  every rung with a constant free of the order already gives the conclusion
  \textup{(}\texttt{MomentLadder.band\_of\_all\_moments}\textup{)}. The dimension restriction on
  the two unconditional bounds is therefore intrinsic to the method, not an artefact of stopping
  early, and the natural a priori input --- the tree walk count --- is false off the coordinate
  case.
\end{itemize}

One question the paper does close, rather than obstruct, is the value of the constant. On the
coordinate locus the mixed characteristic polynomial of a tight rank-$b$ family is the matching
polynomial of an $(a,b)$-biregular incidence graph in the squared variable
(Proposition~\ref{prop:bridge}), so Xu's inequality is a theorem there at every $b$ and, by
Proposition~\ref{prop:xusharp}, the constant $(\sqrt{a-1}+\sqrt{b-1})^2$ is attained and cannot
be lowered. That removes a possibility the numerics had left open --- that the constant is honest
at $b=2$ and slack at $b\ge3$, the adversarial search having stalled at $0.739$ and $0.726$ there
--- and it fixes what any proof of the general statement has to deliver. It also locates the
difficulty: the bridge fails off the coordinate locus, and that is where the whole content of the
conjecture now sits.

What survives all three obstructions is the ratio recursion, which bounds ratios of matching
polynomials rather than Rayleigh quotients, and that is where we would put further effort. The remaining gap is
stated plainly: the certificate closes for six of the seven families tested, and the defect in the
seventh is a right vertex whose child ratio sits well below the interval a uniform certificate can
express. Closing that is the concrete open problem this paper ends on.

\section{Related work}
Two ingredients are classical and are used, not claimed. The factorization of
the characteristic polynomial of a cyclic cover into twisted determinants at
roots of unity is the block-circulant diagonalization
\cite{MS,DFMRS}; in the infinite-cover limit it is the Floquet--Bloch
decomposition of a periodic discrete operator, and $y=-P/(2Q)$ is the
corresponding Floquet discriminant. The identity
$\Phi_{G,r}=\chi_G\cdot\mu_{r-1,G}$ is \cite{HPS}. Closest in form is
McSorley and Feinsilver \cite{MF}. For a path whose edges are cyclically
labelled with $t$ labels they establish a $(\tau,\Delta)$-recurrence
$\Phi_{N}=\tau\Phi_{N-1}-\Delta\Phi_{N-2}$, which is the algebraic skeleton of
\eqref{eq:V}, and their Theorem 4.4 solves it as
$\Delta^{N/2}U_N(\tau/2\sqrt{\Delta})$, exactly the form used here. In the
univariate specialization their $\tau$ is the matching polynomial of the cycle
$C_t$ and $\Delta=(-1)^tx_1\cdots x_t$ is a monomial; passing between their
counting normalization and the characteristic one used here is the
homogenization $V_d\mapsto c^dV_d$, under which a monomial $\Delta$ becomes
$\mu_{G-V(C)}=1$. Their Theorems 4.4 and 4.5 together give
$U_{dn+n-1}=U_d(T_n)U_{n-1}$, the identity invoked in Section~1 to pass between
Theorem~\ref{thm:main} and Hall's conjecture, so that bridge was already
available in 2009. Their parameter is the number of label periods along a path
and their graphs are paths, cycles and rooted trees; there is no covering and
no cycle with trees attached. What is added here is the pair
$(\mu_G,\mu_{G-V(C)})$, whose second entry is a genuine polynomial recording
the attached trees rather than a monomial, together with the covering
interpretation that makes the recurrence a statement about $\mu_{d,G}$.



\section{The formalization}

Everything asserted in this paper is machine-checked in Lean~4/Mathlib unless labelled
otherwise. The sources are \texttt{RamaLean/} in the \texttt{rama-notes} repository; every file
carries its own statement of what is proved and what is assumed, and the standing blocker
throughout is that Mathlib has no matching polynomial and no mixed characteristic polynomial, so
the identities connecting the combinatorial and operator forms are stated rather than derived.
The full file-by-file record, with the axiom base of each theorem, is in \cite{Supp2}.

\begin{thebibliography}{9}
\bibitem{Alek} A.~D.~Aleksandrov, \emph{On the theory of mixed volumes of convex bodies IV.
Mixed discriminants and mixed volumes}, Mat.\ Sb.\ (N.S.) \textbf{3} (1938), 227--251.
\bibitem{Bapat} R.~B.~Bapat, \emph{Mixed discriminants of positive semidefinite matrices},
Linear Algebra Appl.\ \textbf{126} (1989), 107--124.
\bibitem{Gurvits} L.~Gurvits, \emph{The van der Waerden conjecture for mixed discriminants},
Adv.\ Math.\ \textbf{200} (2006), 435--454.
\bibitem{Zhang} J.-C.~Wan, Y.~Wang, Y.-Z.~Fan, \emph{A hypergraph Heilmann--Lieb
theorem}, arXiv:2206.09558 (2022).
\bibitem{XXZ} Z.~Xu, Z.~Xu, Z.~Zhu, \emph{Improved bounds in Weaver's
$\mathrm{KS}_r$ conjecture for high rank positive semidefinite matrices},
arXiv:2106.09205 (2021).
\bibitem{GG} C.~Godsil, I.~Gutman, \emph{On the matching polynomial of a
graph}, in: Algebraic Methods in Graph Theory, 1981.
\bibitem{HPS} C.~Hall, D.~Puder, W.~Sawin, \emph{Ramanujan coverings of
graphs}, Adv.\ Math.\ \textbf{323} (2018), 367--410.
\bibitem{CGHRS} G.~Cochran, C.~Groothuis, A.~Herring, R.~Rohatgi,
E.~Stucky, \emph{A new (combinatorial) proof of the commutativity of
matching polynomials for cycles}, arXiv:1810.05889 (2018). Preprint; no journal
reference or DOI as of August 2026.
\bibitem{WWM} J.~Wan, W.~Wang, A.~Mohammadian, \emph{On the matching
polynomials of graphs with small number of cycles}, arXiv:2103.10580 (2021).
\bibitem{SFM} Y.-M.~Song, Y.-Z.~Fan, Z.~Miao, \emph{Hypergraph coverings and
Ramanujan hypergraphs}, arXiv:2310.01771 (2023).
\bibitem{YY} W.~Yan, Y.-N.~Yeh, \emph{On the matching polynomial of subdivision
graphs}, Discrete Appl.\ Math.\ \textbf{157} (2009) 195--200. The identity used
here is their Corollary 2.1: for $G$ $d$-regular, $m(S(G),x)=x^{m-n}m(G,x^2-d)$.
\bibitem{BDH} G.~Brito, I.~Dumitriu, K.~D.~Harris, \emph{Spectral gap in random
bipartite biregular graphs and applications}, Combin.\ Probab.\ Comput.\
\textbf{31} (2022) 229--267.
\bibitem{HL} O.~J.~Heilmann, E.~H.~Lieb, \emph{Theory of monomer-dimer
systems}, Comm.\ Math.\ Phys.\ \textbf{25} (1972) 190--232.
\bibitem{CC} V.~Y. Chernyak, M.~Chertkov, \emph{Fermions and loops on graphs.
II. A monomer-dimer model as a series of determinants}, J. Stat. Mech. (2008)
P12012.

\bibitem{WC} Y.~Watanabe, M.~Chertkov, \emph{Belief propagation and loop calculus
for the permanent of a non-negative matrix}, J. Phys. A \textbf{43} (2010)
242002.

\bibitem{Je} M.~Jerrum, \emph{Two-dimensional monomer-dimer systems are
computationally intractable}, J. Statist. Phys. \textbf{48} (1987) 121--134.

\bibitem{WF} Y.~Watanabe, K.~Fukumizu, \emph{New graph polynomials from the Bethe
approximation of the Ising partition function}, Combin. Probab. Comput.
\textbf{20} (2011) 299--320.

\bibitem{BC} C.~Bordenave, B.~Collins, \emph{Eigenvalues of random lifts and
polynomials of random permutation matrices}, Ann.\ of Math.\ \textbf{190}
(2019) 811--875.
\bibitem{MF} J.~P.~McSorley, P.~Feinsilver, \emph{Multivariate matching
polynomials of cyclically labelled graphs}, Discrete Math.\ \textbf{309} (2009)
3205--3218.
\bibitem{MS} H.~Mizuno, I.~Sato, \emph{Characteristic polynomials of some graph
coverings}, Discrete Math.\ \textbf{142} (1995) 295--298.
\bibitem{DFMRS} C.~Dalf\'o, M.~A.~Fiol, M.~Miller, J.~Ryan, J.~\v{S}ir\'a\v{n},
\emph{An algebraic approach to lifts of digraphs}, Discrete Appl.\ Math.\
\textbf{269} (2019) 68--76.
\bibitem{MSSII} A.~W. Marcus, D.~A. Spielman, N.~Srivastava,
\emph{Interlacing families II: mixed characteristic polynomials and the Kadison--Singer
problem}, Ann. of Math. (2) \textbf{182} (2015), 327--350.

\bibitem{MSS} A.~Marcus, D.~Spielman, N.~Srivastava, \emph{Interlacing
families I: Bipartite Ramanujan graphs of all degrees}, Ann.\ of Math.\
\textbf{182} (2015), 307--325.
\bibitem{CsF} P.~Csikv\'ari, P.~E.~Frenkel, \emph{Benjamini--Schramm continuity
of root moments of graph polynomials}, European J.\ Combin.\ \textbf{52} (2016),
part B, 302--320.
\bibitem{AH} M.~Ab\'ert, T.~Hubai, \emph{Benjamini--Schramm convergence and the
distribution of chromatic roots for sparse graphs}, Combinatorica \textbf{35}
(2015), 127--151.
\bibitem{Go} C.~D.~Godsil, \emph{Matchings and walks in graphs}, J.\ Graph
Theory \textbf{5} (1981), 285--297.
\bibitem{HallCE} C.~Hall, \emph{A simple bipartite counterexample to Conjecture 10:
  spectral localization of the ordinary matching polynomial}, verification note, personal
  communication, August 2026.
\bibitem{Avakov} E.~R.~Avakov, \emph{Extremum conditions for smooth problems with
equality-type constraints}, USSR Comput.\ Math.\ Math.\ Phys.\ \textbf{25} (1985), 24--32.

\bibitem{Arutyunov} A.~V.~Arutyunov, \emph{Optimality Conditions: Abnormal and Degenerate
Problems}, Kluwer, 2000.

\bibitem{RL} M.~Ravichandran, J.~Leake, \emph{Mixed determinants and the Kadison--Singer
  problem}, Math. Ann. \textbf{377} (2020), 511--541.
\bibitem{AFH} O.~Angel, J.~Friedman, S.~Hoory, \emph{The non-backtracking spectrum of the
  universal cover of a graph}, Trans.\ Amer.\ Math.\ Soc.\ \textbf{367} (2015), 4287--4318;
  arXiv:0712.0192.
\bibitem{BCs} F.~Bencs, P.~Csikv\'ari, \emph{Evaluations of Tutte polynomials of
regular graphs}, arXiv:2105.06798 (2021).
\bibitem{Supp2} V.~Bonfioli, \emph{The biregular case: supplementary material},
  companion note, 2026.
\bibitem{PartI} V.~Bonfioli, \emph{Roots of $d$-matching polynomials and the spectrum of
  the universal cover: a closed form at first Betti number one, and a refuted conjecture},
  companion paper, 2026.
\end{thebibliography}

\end{document}
